\documentclass[11pt]{article}
\usepackage[a4paper,margin=28mm]{geometry}
\usepackage{amsmath,amssymb,amsthm,mathtools,booktabs}
\usepackage[T1]{fontenc}
\usepackage{lmodern}
\DeclareMathSizes{10}{10}{8}{8}
\DeclareMathSizes{10.95}{10.95}{8}{8}
\usepackage[hidelinks]{hyperref}
\usepackage{microtype}
\newtheorem{theorem}{Theorem}[section]
\newtheorem{lemma}[theorem]{Lemma}
\newtheorem{proposition}[theorem]{Proposition}
\newtheorem{corollary}[theorem]{Corollary}
\theoremstyle{remark}\newtheorem{remark}[theorem]{Remark}
\newcommand{\wt}[1]{\lVert #1\rVert}
\newcommand{\len}{\ell}
\newcommand{\coef}[2]{[z^{#1}]#2}
\newcommand{\Del}{\Delta}
\newcommand{\Iv}[2]{I_{[#1,#2]}}
\newcommand{\rev}{^{\!*}}
\newcommand{\RR}{\mathcal R}
\DeclareMathOperator{\supp}{supp}
\allowdisplaybreaks
\title{Effective rank-unimodality under dilation and\\cumulative inequalities for composition ideals}
\author{}
\date{Research manuscript, 15 October 2026}
\begin{document}
\maketitle
\begin{abstract}
The rank polynomial of a composition counts distinct smaller compositions in generalized subword order, by total weight. We prove that every fixed composition becomes rank-unimodal after sufficiently large integer dilation. The threshold is computable by exact rational arithmetic. For a composition of weight $W$ and length $l$ whose largest part is less than $W/2$, every sufficiently large dilation has its mode at $\lfloor(tW+l-1)/2\rfloor$, with at most a tie at the next integer; the eventual choice is determined exactly in each parity. We also prove log-concavity of tail counts for all one-descent and bitonic compositions, and for arbitrarily many valleys separated by a common maximum. A sharp growth-deficit inequality holds for every composition downset missing a word of a specified weight. For the rank coefficients themselves, arbitrary hosts are unimodal when $N\le2M+5$, where $M$ is the largest part. Finally, strict one-descent cores of weight less than $160M^3$ suffice to verify rank-unimodality at a fixed descent peak $M$, and exact finite profiles settle all peaks at most seven. Analytic proofs, reproducible certificates, and selected unbounded Lean proofs accompany the results.
\end{abstract}

\section{Statement and scope}
A composition is a finite word of positive integers; an empty maximum is taken to be zero. Write $u\preceq w$ if the letters of $u$ can be embedded, in order, into distinct positions of $w$, with each letter no larger than its host letter. Set
\[
 F_w(z)=\sum_{u\preceq w}z^{\wt u},\qquad F_\varnothing=1,
 \qquad S_t=z+z^2+\cdots+z^t.
\]
Every word $u$ occurs once in this sum, regardless of its number of embeddings. A polynomial is unimodal if its coefficient sequence is weakly increasing and then weakly decreasing. A word is bitonic if its parts weakly increase to a maximum and then weakly decrease. Coefficients throughout are extended by zero to all integer indices. Write $\Del P(n)=P_n-P_{n-1}$.

\begin{theorem}[Effective dilation]\label{thm:dilationmain}
For every fixed nonempty composition $w=(w_1,\ldots,w_l)$ there is a
computable $T(w)$ such that $F_{(tw_1,\ldots,tw_l)}$ is unimodal for
every integer $t\ge T(w)$.
If $W=\wt w$ and $2\max_iw_i<W$, the mode is
$\lfloor(tW+l-1)/2\rfloor$, possibly tied with
$\lceil(tW+l-1)/2\rceil$. The eventual alternative in each parity of
$t$ is computable from explicit rational polynomials.
\end{theorem}

The mode assertion also holds for $t\ge T(w)$, after enlarging the
same effective threshold if necessary. Section~\ref{sec:dilation}
constructs a threshold for both claims simultaneously. Its proof uses
a universal right-tail bound and a central window whose width depends
only on $l$. In particular, it establishes eventual unimodality for
every shape, with arbitrarily many descents.

\begin{theorem}[Cumulative rank regularity]\label{thm:cumulativemain}
If $w$ has at most one adjacent descent, or is bitonic, then its tail counts
$t_k=\sum_{j\ge k}[z^j]F_w$ satisfy $t_k^2\ge t_{k-1}t_{k+1}$ for
every integer $k$. There is no restriction on the length or part sizes.
The same conclusion holds when $w=Z_0MZ_1M\cdots MZ_s$, where every
$Z_i$ weakly decreases and then weakly increases, with parts below $M$; in that
case $1+(S_M-1)F_w$ also has log-concave tail counts.
\end{theorem}

\begin{theorem}[Sharp rank-growth deficit]\label{thm:deficitmain}
Let $f_k$ count weight-$k$ compositions in a downset, and suppose a
composition of weight $\kappa\ge2$ is missing. For every $k\ge\kappa-1$,
\[
 (2^{\kappa-1}-1)(2f_k-f_{k+1})\ge f_k.
\]
The constant is best possible, even among finite principal ideals.
Moreover, for $k\ge1$, equality $f_{k+1}=2f_k>0$ forces both ranks
to be full: $f_k=2^{k-1}$ and $f_{k+1}=2^k$.
\end{theorem}

Theorems~\ref{thm:cumulativemain} and~\ref{thm:deficitmain} are proved in
Sections~\ref{sec:cumulativecores} and~\ref{sec:saturation}. Tail-count
log-concavity does not imply rank-unimodality: for example, the
non-unimodal sequence $(10,10,8,9,1)$ has log-concave tail sums
$(38,28,18,10,1)$. The next three results concern the rank coefficients themselves.

\begin{theorem}[Dominant maximum]\label{thm:dominantmain}
An arbitrary composition of weight $N$ and maximum part $M$ has a
unimodal rank polynomial whenever $N\le2M+5$.
\end{theorem}

This result places no restriction on the number of descents. Its
analytic proof in Section~\ref{sec:dominant} compares the ranks at
distance two from the top with the bottom three ranks.

\begin{theorem}\label{thm:main}
Let $w$ have at most one adjacent descent. If it has a descent $w_i>w_{i+1}$, assume $w_i\le7$. Then $F_w$ is unimodal. In particular every one-descent composition with all parts at most seven is rank-unimodal. Every bitonic composition with maximum part at most seven is also rank-unimodal.
\end{theorem}

\begin{theorem}[Finite verification principle]\label{thm:finiteprinciple}
Fix $M\ge4$. Suppose $F_{X,M,Y}$ is unimodal for all increasing arms $X,Y$ with parts strictly below $M$ and $\wt X+M+\wt Y<160M^3$. Then every one-descent composition of descent peak $M$ has a unimodal rank polynomial.
\end{theorem}

Theorem~\ref{thm:finiteprinciple} follows from the stronger relative statement in Theorem~\ref{thm:cubicbound}. Its weight threshold is polynomial in $M$, but the number of cores to be checked need not be. The theorem does not assert that these finite checks have been completed for every peak.

The bounded-peak claim concerns arbitrary lengths, not merely compositions with seven parts. The dilation theorem covers every fixed shape but does not assert that dilation is unnecessary. None of the results resolves the unrestricted composition conjecture or asserts tail log-concavity of the rank coefficients. The dilation, dominant-maximum, and deficit proofs are analytic. The bounded-peak result additionally uses the supplied finite arithmetic certificates. The cumulative theorem invokes the discrete cumulative-convolution closure theorem of Nanda and Sengupta~\cite{NandaSengupta}.

Rank-unimodality of arbitrary composition ideals is the layered-permutation case of a broader pattern-poset question; see Vatter~\cite{Vatter}, especially Conjecture~2.6. Ordinary unweighted subsequence counts are log-concave by Chase's theorem; Vatter~\cite{VatterWords} gives a short proof. Weight grading behaves differently: the initial coefficients $1,1,2$ already rule out coefficient log-concavity. Monotone composition hosts are known to be unimodal. For rectangular hosts this is Sagan's composition-in-a-rectangle theorem~\cite{Sagan}; his more general ideal theorem uses row-wise Ferrers containment, whereas the order here allows host positions to be skipped.

Zhang~\cite{Zhang} gives the symmetric maximum decomposition and several one-descent families. Chen's public abstract~\cite{Chen} reports unimodality for arbitrary compositions with at most six parts, tail log-concavity for bitonic hosts with parts at most five, and an $O(M2^M)$ weight bound for a least counterexample to a relative tail-log-concavity assertion. Our comparison with Chen is restricted to that abstract. A matching anonymized sixteen-page version of Zhang's manuscript was inspected; its identity with the complete public version has not been independently confirmed. These source limits are recorded in the bibliography and accompanying provenance notes; no claim about either inaccessible complete version is made.

The dilation argument covers every fixed shape and determines its eventual mode. The cumulative inequalities cover the stated structural classes without a peak bound. The finite-profile argument in Sections~\ref{sec:profiles}--\ref{sec:central} closes peaks six and seven, while Section~\ref{sec:cubic} bounds the entire minimal relative-unimodality obstruction cubically. This cubic bound is not a bound for the stronger relative tail-log-concavity assertion in Chen's abstract. Maximum decompositions, interval layers, and the existing peeling lemmas are background tools, with the needed proofs included. The unrestricted undilated conjecture remains open.

\section{Counting identities and interval layers}\label{sec:basic}
\begin{lemma}[Rightmost last letter]\label{lem:last}
For a nonempty host $w$, let $r_t$ be the last position with $w_{r_t}\ge t$. Then
\[
 F_w=1+\sum_{t=1}^{\max w}z^tF_{w_1\cdots w_{r_t-1}}.
\]
In particular, appending a global maximum $t$ gives $F_{wt}=1+S_tF_w$.
\end{lemma}
\begin{proof}
Every nonempty smaller word has a unique last letter $t$. Move its last embedded letter to position $r_t$; its preceding letters still embed in the indicated prefix. This is a characterization of distinct words, not an enumeration of embeddings.
\end{proof}

\begin{lemma}[Maximum decomposition]\label{lem:core}
For arbitrary words $X,Y$ and $M\ge\max(X,Y)$,
\begin{equation}\label{eq:core}
 F_{X,M,Y}=F_X+F_Y+(S_M-1)F_XF_Y.
\end{equation}
\end{lemma}
\begin{proof}
Put $A=F_X$ and $K=1+(S_M-1)A$. We prove $F_{X,M,Y}=A+KF_Y$ by induction on the length of $Y$, simultaneously for all its prefixes. The empty case is $1+S_MA=A+K$. On appending a letter $v\le M$, use Lemma~\ref{lem:last}. The contributions involving $A$ sum to $S_MA$, while those involving $K$ sum to $K(F_{Yv}-1)$. Thus the result is $1+S_MA+K(F_{Yv}-1)=A+KF_{Yv}$.
\end{proof}

For weakly increasing arms $X,Y$, all their parts strictly below $M$, use the notation
\begin{equation}\label{eq:notation}
 a=\wt X,\quad b=\wt Y,\quad D_Y=1+(S_M-1)F_Y,
 \quad Q_{X,Y}=F_XD_Y,\quad H=F_Y+Q_{X,Y}.
\end{equation}
Then $H$ is symmetric in the arms. If $X=Zu$, the two append recurrences are
\begin{equation}\label{eq:qappend}
 Q_{Zu,Y}=D_Y+S_uQ_{Z,Y},\qquad
 Q_{X,Yv}=S_vQ_{X,Y}+(S_M-S_v)F_X.
\end{equation}

For integers $l\le r$, let $\Iv l r=z^l+\cdots+z^r$. A nonnegative integer unimodal polynomial has a height decomposition
\[
 P=\sum_{h=1}^{\max P_i}\Iv{l_h}{r_h},\qquad
 \{i:P_i\ge h\}=[l_h,r_h]\cap\mathbb Z.
\]
The convolution of two intervals is symmetric and unimodal about the sum of their centers. The same holds for a product of any number of intervals. Consequently, if all layer endpoint sums of $P$ and $Q$ are at least $c$ and $d$, then $PQ$ is nondecreasing through $\lfloor(c+d+1)/2\rfloor$. This does not imply unimodality of the entire product.

We repeatedly use the following elementary equivalence: for an already unimodal nonnegative polynomial $P$, all its layer endpoint sums are at least $c$ if and only if
\begin{equation}\label{eq:comparison}
 P_i\le P_j\quad\text{whenever }i\le j\text{ and }i+j\le c.
\end{equation}
Indeed, each layer interval containing $i$ then contains $j$. Conversely, if $l+r<c$ for a layer, its left endpoint is present while its reflection $c-l>r$ is absent.

\begin{lemma}[Monotone arms]\label{lem:arms}
For a weakly increasing arm $Z$ of weight $c$, $F_Z$ is unimodal, its lowest layer is $[0,c]$, and every other nonempty layer has endpoint sum at least $c+1$. If $F_Z$ is not a single interval, its second layer is exactly $[2,c-1]$.
\end{lemma}
\begin{proof}
Use $F_{Zt}=1+S_tF_Z$. Multiplication by an interval preserves unimodality: its differences are $p_n-p_{n-t}$, and these cannot change from negative to positive when $p$ is unimodal. The product has first and last nonzero coefficient one. Adding the constant term one therefore preserves unimodality. By induction, its symmetric interval-product summands have centers at least $(c+t+1)/2$; this gives the asserted reflection for the higher layers. The lowest layer becomes $[0,c+t]$. The polynomial is not an interval exactly when the host has at least two parts and at least one part greater than one. In that case the coefficients at ranks two and $c-1$ are at least two, whereas ranks zero, one and $c$ have coefficient one. Unimodality identifies the second layer.
\end{proof}

\begin{lemma}[First-half increase]\label{lem:half}
For an arbitrary host of weight $N$, its rank polynomial is nondecreasing through rank $\lceil N/2\rceil$.
\end{lemma}
\begin{proof}
Replace the host by the layered permutation with decreasing block sizes given by its parts. Its distinct induced patterns are exactly the smaller compositions, under this correspondence. On real functions on the $k$-subsets of its $N$ positions, let $U$ sum over one-element deletions and let $D$ be its adjoint. On rank $k$ one has $DU-UD=(N-2k)I$. Hence $U$ is injective when $k<N/2$, since $\|Uf\|^2=\|Df\|^2+(N-2k)\|f\|^2$. The subspace of functions constant on subsets inducing the same pattern is carried by $U$ into the analogous next subspace: deletion counts depend only on the induced pattern. Its dimension is the number of distinct patterns of that rank. Restricting the injection proves the assertion. This is the usual first-half argument discussed in~\cite{Vatter}.
\end{proof}

\section{Effective unimodality under dilation}\label{sec:dilation}

For a fixed nonempty composition $w=(w_1,\ldots,w_l)$, write
$tw=(tw_1,\ldots,tw_l)$ and $W=\sum_iw_i$.

\begin{theorem}[Dilation and its eventual modes]\label{thm:dilation}
There is a computable integer $T(w)$ such that $F_{tw}$ is unimodal
for every integer $t\ge T(w)$.
If $2\max_iw_i<W$, set $C_t=(tW+l-1)/2$. For all $t\ge T(w)$,
the lower integer $\lfloor C_t\rfloor$ is a mode, and it is either
the unique mode or belongs to exactly two adjacent modes
$\lfloor C_t\rfloor,\lceil C_t\rceil$.
In particular, an integral $C_t$ is the unique mode.
The eventual choice in each parity of $t$ is computable exactly.
There is no restriction on the length or number of descents of $w$.
\end{theorem}

The proof confines all undecided differences to a fixed finite window,
then computes their exact polynomials in $t$. A limiting unimodal
density, without control of these differences, would not suffice.

\begin{lemma}[A universal right-tail bound]\label{lem:righttail}
For a nonempty host of weight $N$ and length $l$, its rank polynomial
is a nonnegative sum of products of interval polynomials, each
symmetric about a point at most $(N+l-1)/2$.
Consequently
\[
 [z^n]F_w\le[z^{n-1}]F_w
       \qquad\left(n\ge\left\lceil\frac{N+l}{2}\right\rceil\right).
\]
\end{lemma}
\begin{proof}
For $l=1$, use $F_w=\Iv0N$. Assume $l\ge2$. Set
$m_i=\max(w_{i+1},\ldots,w_l)$, with $m_l=0$, and
$N_i=w_1+\cdots+w_i$. Group the terms of the rightmost-last-letter
identity according to their host position:
\[
 F_w=1+\sum_{i:w_i>m_i}\Iv{m_i+1}{w_i}F_{w_1\cdots w_{i-1}}.
\]
For $i\ge2$, the induction hypothesis bounds twice the center of
each product in this summand by
\[
 N_{i-1}+i-2+(w_i+m_i+1)=N_i+i-1+m_i\le N+l-1,
\]
since $m_i\le N-N_i$. For $i=1$, the twice-center is
$w_1+m_1+1\le N+l-1$, using $l\ge2$.
The constant term also meets the bound. Products of intervals are
symmetric and unimodal, so each is nonincreasing at every difference
index $n\ge\lceil(N+l)/2\rceil$. Summing proves the assertion.
\end{proof}


Together with Lemma~\ref{lem:half}, the right-tail bound leaves only
the differences strictly after $\lceil N/2\rceil$ and no later than
$\lfloor(N+l-1)/2\rfloor$ undecided. The width of this window depends
on the number of parts, not their sizes.

\begin{lemma}[Multilinear expansion]\label{lem:multilinear}
There are integers $a_I$, independent of $t$, such that
\[
 F_{tw}=\sum_{I\subseteq[l]}a_I\prod_{i\in I}S_{tw_i},
 \qquad a_{[l]}=1,\qquad
 \sum_{|I|=j}a_I=1\quad(0\le j\le l).
\]
One may choose them with $\sum_I|a_I|\le3^l$.
\end{lemma}
\begin{proof}
Fix recursive splits at maxima of $w$, which remain valid under
dilation. Apply the maximum identity, retaining a formal variable for
each position. The resulting polynomial is multilinear and its full
monomial has coefficient one. Setting every variable equal to $s$
gives $1+s+\cdots+s^l$: for arms of lengths $a,b$, the expressions
$1+(s-1)P_X=s^{a+1}$ and $1+(s-1)P_Y=s^{b+1}$ multiply.
This proves the sums by cardinality.
An endpoint maximum gives $P=1+sP_Y$. At an internal maximum,
$\|P\|_1\le\|P_X\|_1+\|P_Y\|_1+2\|P_X\|_1\|P_Y\|_1$.
For $a,b\ge1$, $3^a+3^b\le3^{a+b}$, so induction gives the norm bound.
\end{proof}

For positive widths $v_1,\ldots,v_m$, let
$B_v=\mathbf1_{[0,v_1]}*\cdots*\mathbf1_{[0,v_m]}$.
For $m\ge2$ its continuous representative is
\begin{equation}\label{eq:boxspline}
 B_v(x)=\frac1{(m-1)!}
 \sum_{J\subseteq[m]}(-1)^{|J|}
       \left(x-\sum_{j\in J}v_j\right)_+^{m-1}.
\end{equation}
Repeated integration of interval indicators proves this formula.

\begin{lemma}[Negative central curvature]\label{lem:boxcurvature}
If $m\ge3$ and $2\max_i v_i<V=\sum_i v_i$, then
$B_v''(V/2)$ exists and is strictly negative.
\end{lemma}
\begin{proof}
We use the elementary strict shape of interval convolutions. In
centered coordinates, a convolution with total width $V$ and largest
width $L$ has a plateau of half-width $(2L-V)_+/2$ and is strictly
decreasing from that plateau to its support endpoint. For two
intervals this is the trapezoid formula. Inductively put the largest
interval last, and write $g$ for the centered convolution of the other
intervals, of total width $D=V-L$. Then
\[
 B_v'(V/2+x)=g(x+L/2)-g(x-L/2).
\]
For $(L-D)_+/2<x<(L+D)/2$, the inner point lies inside the support,
and the outer point lies beyond the plateau of $g$. That plateau has
half-width strictly less than $L/2$ when at least two widths remain.
Symmetry and the inductive strict decrease make the displayed
difference negative. At $x=0$ the difference is zero by symmetry.
If $L\ge D$, both terms vanish for $0\le x\le(L-D)/2$.
This proves the asserted strict shape, including strict negativity
of its derivative off the plateau for three or more intervals.

Under $L<D$, the point $L/2$ lies strictly between the plateau and
the support endpoint of $g$. Hence
\[
 B_v''(V/2)=2g'(L/2)<0.
\]
When $g$ is a two-interval convolution the point is in a sloping
open segment of its trapezoid; otherwise $g'$ is continuous there.
\end{proof}

Write $c_i=a_{[l]\setminus\{i\}}$, so $\sum_i c_i=1$.
The top two multilinear degrees have an exact symmetry:
\begin{equation}\label{eq:topsymmetry}
 H_t=\prod_i S_{tw_i}+\sum_i c_i\prod_{j\ne i}S_{tw_j},
 \qquad z^{tW+l-1}H_t(1/z)=H_t(z).
\end{equation}
Indeed, reflection multiplies the full product by $z^{-1}$ and the
$i$th omitted-factor product by $z^{tw_i}$. The differences cancel,
using $z^{tw_i}-1=-(1-z)S_{tw_i}/z$ and $\sum_i c_i=1$.

\begin{lemma}[Exact central polynomials]\label{lem:centralpolynomials}
Assume $2\max_iw_i<W$. Choose a parity $\epsilon\in\{0,1\}$ and
a half-integer $s$ with $0\le s\le l/2+1$ and
$\epsilon W/2+s\in\mathbb Z$. There is an explicitly computable
rational polynomial $P_s(t)$ such that
\[
 \Del F_{tw}(tW/2+s)=P_s(t)
 \quad\text{for }t\ge4l+4,\quad t\equiv\epsilon\pmod2.
\]
For $l\ge4$, its degree is at most $l-3$, and its coefficient of
$t^{l-3}$ is
\begin{equation}\label{eq:centralleading}
 B_w''(W/2)(s-l/2).
\end{equation}
For $l=3$ the polynomial is exactly $-2s+3$.
At $s=l/2$, when this parity exists, write $Q_w=P_{l/2}$.
This polynomial is either zero or has negative leading coefficient;
for $l\ge4$ its degree is at most $l-4$.
\end{lemma}
\begin{proof}
For a product of $m\ge2$ intervals the exact first-difference formula is
\begin{equation}\label{eq:centralbinomial}
 \Del\left(\prod_{i\in I}S_{tw_i}\right)(n)
 =\sum_{J\subseteq I}(-1)^{|J|}
    \binom{n-t\sum_{j\in J}w_j-2}{m-2}.
\end{equation}
A binomial is zero unless its upper argument is at least its
nonnegative lower argument. This follows by expanding the numerator
of $z^m\prod_{i\in I}(1-z^{tw_i})/(1-z)^{m-1}$.
At $n=tW/2+s$, a summand has upper argument
\[
 t\beta+s-2,\qquad \beta=W/2-\sum_{j\in J}w_j.
\]
For $t\ge4l+4$, a positive $\beta$ is at least $1/2$ and makes the
binomial an ordinary polynomial in $t$; a negative $\beta$ is at
most $-1/2$ and makes it zero. If $\beta=0$, retain its constant
binomial if it is nonzero. In that case $W$ is even and the allowed
parity makes $s$ integral. Products with fewer than two factors have
zero difference in this balanced central window. This constructs
$P_s$ exactly from Lemma~\ref{lem:multilinear}.

For $l\ge4$, the leading two terms of the full product's difference
are
\[
 t^{l-2}B_w'(W/2)
 +t^{l-3}\left(s-\frac{l+1}{2}\right)B_w''(W/2).
\]
To verify the second coefficient directly, the shifts in the
binomial product are $2,3,\ldots,l-1$, whose sum is
$(l-2)(l+1)/2$. The first term vanishes by symmetry.
An omitted-factor product contributes leading coefficient
$B_{w\setminus i}'(W/2)$ at degree $l-3$. Differentiating the
convolution with the $i$th interval and using symmetry gives
\[
 B_w''(W/2)=B_{w\setminus i}'(W/2)
              -B_{w\setminus i}'(W/2-w_i)
           =2B_{w\setminus i}'(W/2).
\]
All other products are of degree at most $l-4$ after taking a
difference. A zero-slope binomial is only constant and cannot affect
the coefficient at degree $l-3$. Summing, with $\sum_i c_i=1$,
proves~\eqref{eq:centralleading}.

For $l=3$, all singleton subset sums are below $W/2$ and all pair
sums are above it. The full product contributes $-2(s-2)$, each
two-factor product has difference $-1$, and their coefficients sum
to one. This gives $-2s+3$ exactly.

Finally, the degree bound for $Q_w$ follows either from
\eqref{eq:centralleading} or the exact cancellation
\eqref{eq:topsymmetry}. Lemma~\ref{lem:righttail} gives $Q_w(t)\le0$
at all integers of the indicated parity where it represents the
central difference. A nonzero polynomial with this property has
negative leading coefficient.
\end{proof}

For a nonzero rational polynomial $P(t)=\sum_{j=0}^d p_jt^j$, define
\[
 b(P)=1+\left\lceil\frac{\sum_{j<d}|p_j|}{|p_d|}\right\rceil,
 \qquad b(0)=1.
\]
For every real $t\ge b(P)$, a nonzero $P(t)$ has the sign of $p_d$.
Indeed, if $d\ge1$ and $t\ge1$, the absolute sum of its lower terms
is at most $t^{d-1}\sum_{j<d}|p_j|<|p_d|t^d$.
The constant case is immediate.

\begin{proof}[Proof of Theorem~\ref{thm:dilation}, including an effective threshold]
If $2\max_iw_i\ge W$, split a dilated word at its maximum $M$ and
write $D$ for the remaining weight. Here $M\ge D$. First-half
increase reaches rank $D$, whereas the maximum identity gives
\[
 \Del F_{XMY}(n)=-[z^{n-M-1}]F_XF_Y\le0
 \qquad(n\ge D+2).
\]
Only the difference at $D+1$ is undecided, so take $T(w)=1$.
Otherwise $l\ge3$ and the central curvature is negative.

Fix a parity of $t$, and write $N=tW$, $L=\lceil N/2\rceil$,
$V=\lfloor(N+l-1)/2\rfloor$, $U=\lceil(N+l-1)/2\rceil$.
The offsets $n-N/2$ in the following list depend only on the parity:
\[
 L<n\le V;\qquad n=V;\qquad n=U+1.
\]
For the first two groups, $s<l/2$, so $P_s$ has positive leading
coefficient. For the last index, $s>l/2$, so the leading coefficient
is negative. Include $Q_w$ as well when $N+l$ is even.
Take $T(w)$ to be the maximum of $4l+4$ and all the integers $b(P)$
for these finitely many polynomials, over both parities.
This is a finite exact-arithmetic construction.

For $t\ge T(w)$, Lemma~\ref{lem:half} gives nonnegative differences
through $L$; the selected positive polynomials give positive
differences from $L+1$ through $V$, as well as at $V$ itself.
Lemma~\ref{lem:righttail} gives nonpositive differences at every
index after $V$. Hence the sequence is unimodal and $V$ is a mode.
The strict difference at $V$ excludes earlier modes, and the strict
negative difference at $U+1$ excludes later modes.
If $U=V$, this proves uniqueness. Otherwise the difference at $U$
is $Q_w(t)$; it is either strictly negative beyond its bound or
identically zero. This proves the precise eventual alternatives.
\end{proof}

For example, the supplied exact procedure gives $T(w)=32$ for
$w=(1,4,2,5,3,2,4)$, with unique eventual mode
$\lfloor(21t+6)/2\rfloor$ in both parities. This is a sufficient
threshold, not a claim of the least possible threshold.
The procedure computes and certifies central polynomials; it does not
test finitely many dilations and extrapolate.

\begin{corollary}[Mean rank]\label{cor:dilationmean}
Choose a distinct smaller composition uniformly from the ideal of
$tw$, including the empty composition, and let its weight be $K_t$.
For every fixed nonempty $w$,
\[
 \mathbb E K_t=\frac{tW+l-1}{2}+O(t^{-1}).
\]
\end{corollary}
\begin{proof}
A product $\prod_{i\in I}S_{tw_i}$ has mass
$t^{|I|}\prod_{i\in I}w_i$ and mean exponent
$(t\sum_{i\in I}w_i+|I|)/2$. The full polynomial has mass
$t^l\prod_iw_i+O(t^{l-1})$.
Subtract $(tW+l)/2$ times this mass from its first moment.
The terms missing exactly one factor contribute
\[
 -\tfrac12t^l\prod_iw_i\sum_i c_i+O(t^{l-1})
 =-\tfrac12t^l\prod_iw_i+O(t^{l-1});
\]
all other proper terms contribute $O(t^{l-1})$.
Divide by the total mass. The same calculation holds for $l=1$,
when the sole proper term is the empty product.
\end{proof}


\section{The relative kernel}\label{sec:kernel}
We include the kernel estimates needed later, so that no unproved shape-closure principle is being used.

\begin{lemma}\label{lem:kernel}
For a weakly increasing arm $Y$ with parts below $M$, put $b=\wt Y$ and $R=b+M$. The polynomial $D_Y$ is nonnegative and unimodal, of degree $R$. Every height layer has endpoint sum at least $R+1$. Thus
\begin{equation}\label{eq:qgrowth}
 Q_{X,Y}\text{ is nondecreasing through }
 \left\lfloor\frac{a+b+M+2}{2}\right\rfloor.
\end{equation}
For nonempty $Y$, $d_0=d_1=0$, $d_R=1$, and the positive support of $D_Y$ is an interval.
\end{lemma}
\begin{proof}
For an arm $Z$ define $E_Z=1-z+(2z-1)F_Z$. Then
\[
 E_\varnothing=z,\quad E_{Zt}=S_tE_Z+z^{t+1},\qquad
 (1-z)D_Z=E_Z-z^{M+1}F_Z.
\]
The positive support of $E_Z$, for weight $c$, is an interval $[q_Z,c+1]$, with $q_Z\le\max Z+1$. This follows directly by induction from the displayed recurrence. In particular $D_Z$ is nondecreasing through rank $M$.

Write $Y=Zt$, $c=\wt Z$, $P=D_Z$, and
\[
 D_Y=S_tP+J,\qquad J=S_M-S_t.
\]
This proves nonnegativity inductively. The only possible splice after the already increasing part is at $M+1$, with difference $P_M-P_{M-t}-1$. If $P_M-P_{M-t}\ge1$, the splice is increasing and the subsequent coefficients agree with the unimodal $S_tP$. Otherwise $P_M=P_{M-t}$, and the interval support of $E_Z$ forces $c+1\le M-t$. The polynomial $P$ is nonincreasing thereafter, so every subsequent window difference $P_{n-1}-P_{n-t-1}$ is nonpositive. This proves unimodality, including the zero-plateau case.

For the reflection, induct on the length of $Y$; $D_\varnothing=S_M$ is the base. Set $d=c+M$, and assume $P_j\le P_{d+1-j}$ on the left half. Each interval-product summand of $W=S_tP$ has center at least $(d+t+2)/2$, so it already obeys reflection about $(d+t+1)/2$. Only a unit loss from $J$ remains. Such a loss can occur only when
\[
 t+1\le i\le M,\quad i\le c+t,\quad 2i<d+t+1.
\]
Write $K_j=P_{d+1-j}-P_j$. Direct summation gives
\[
 W_{d+t+1-i}-W_i=\sum_{j=i-t+1}^{i}K_j+P_i-P_{i-t}.
\]
The window endpoint sum is at most $d+1$, so its negative $K_j$ terms pair with their nonnegative reflected terms inside the window. The last difference equals a sum of coefficients of $E_Z$ over $[i-t+1,i]$. This interval meets $[q_Z,c+1]$, by the displayed conditions, so that difference is at least one. It pays for the possible unit loss. This proves the layer bound. Combining it with Lemma~\ref{lem:arms} proves~\eqref{eq:qgrowth}.
\end{proof}

\begin{lemma}[Strict tail budget]\label{lem:budget}
If $F_Y=\sum p_jz^j$, then for $1\le t\le\min(b-1,M-1)$,
\begin{equation}\label{eq:budget}
 d_{R-t}-d_{t+1}=\sum_{j=0}^t(p_{b-j}-p_j)+p_{t+1}\ge p_{t+1}\ge1.
\end{equation}
\end{lemma}
\begin{proof}
The equality follows from the definition of $D_Y$. Ordinary reflection $p_j\le p_{b-j}$ holds on the left half by Lemma~\ref{lem:arms}. In the prefix sum past the midpoint, cancel reflected pairs; the remaining terms are nonnegative. All ranks between zero and $b$ occur.
\end{proof}

\begin{lemma}[Two stronger reflections]\label{lem:strongkernel}
If $b\ge M$, every layer of $D_Y$ has endpoint sum at least $R+2$. If $b\ge2M-1$ and $Y$ has at least two non-one parts, the endpoint sums are at least $R+3$.
\end{lemma}
\begin{proof}
Write $Y=Zt$, $c=b-t$, $r=c+M$, and $E=D_Z$. For the first claim $c\ge1$. The product $W=S_tE$ already has reflection sum $R+2$. Only $i\in[t+1,M]$ needs a unit payment for $J=S_M-S_t$; its reflected index is beyond $M$. Put
\[
 f_j=e_{r+1-j}-e_j,\quad L=i-t,\quad U=i-1.
\]
The reflection difference of $W$ is $\sum_{j=L}^Uf_j$. Here $1\le L\le c$ and $L+U\le r-1$. Pair all negative terms with their reflections. The term $f_L$ remains unpaired. If $L=1$, it is $e_r-e_1=1$; otherwise Lemma~\ref{lem:budget} gives $f_L\ge1$. This proves the first claim.

For the second claim $t\ge2$, $c\ge M$, and $F_Z$ is not an interval. By the first claim $W$ has reflection sum $R+3$. Put $g_j=e_{r+2-j}-e_j$. Then $g_1=0$, while, for $2\le j\le M-1$,
\[
 g_j=\sum_{h=0}^{j-2}p_{c-h}-\sum_{h=0}^{j-1}p_h+p_j
 \ge p_j-1\ge1.
\]
The first inequality is the cumulative strengthened reflection of Lemma~\ref{lem:arms}; the second uses its second layer. At a correction index $i\in[t+1,M]$, the reflection difference is $\sum_{j=i-t}^{i-1}g_j$. It contains at least two terms, all in $[1,M-1]$, so is at least one. This completes the payment and the proof.
\end{proof}

\section{Exact saturation and sharp growth deficits}\label{sec:saturation}

The following statement applies to any finite downset of compositions,
not only to a principal ideal or a monotone host. Let $f_k$ count its
members of weight $k$.

\begin{theorem}[Saturation rigidity at every rank]
Let $k\ge1$ and assume $f_k>0$. If $f_{k+1}=2f_k$, then the downset
contains every composition of weights $k$ and $k+1$. In particular,
$f_k=2^{k-1}$ and $f_{k+1}=2^k$.
\end{theorem}

\begin{proof}
Decrement the last part of a word of weight $k+1$, deleting it if it is
one. Every word $v$ of weight $k$ has at most two preimages: append a
one, or increase its last part. Equality of the total counts forces
both preimages of every feasible $v$ to be feasible.

Now remove one unit from the first part of either preimage, deleting
that part if necessary. Downward closure keeps the resulting word
feasible. Encode a composition of weight $k$ by its $k-1$ cut bits.
These two operations delete the first cut bit and append, respectively,
a one or a zero. Starting from any feasible cut string, $k-1$ successive
operations, with the appended bits prescribed, produce any desired cut
string of length $k-1$. Thus every composition of weight $k$ is feasible.
For $k=1$ the cut string is empty and the same conclusion is immediate.
The two already forced extensions then give every composition of
weight $k+1$.
\end{proof}

\begin{corollary}[The exact support of the growth deficit]
For a nonempty host $w$ of weight $N$, let $\kappa(w)$ be the smallest
weight of a composition not below $w$. Put
\[
 E_w=1-z+(2z-1)F_w.
\]
Then the positive coefficient support of $E_w$ is precisely
$[\kappa(w),N+1]\cap\mathbb Z$. Consequently
\[
 (F_w)_{n+1}\le2(F_w)_n-1
 \quad\text{for }\kappa(w)-1\le n\le N.
\]
Moreover $2\le\kappa(w)\le\min\{\len w+1,\max w+1\}$.
\end{corollary}

\begin{proof}
All ranks through $N$ are nonempty, and the two-preimage map gives
$f_{j}\le2f_{j-1}$ for $j\ge2$. Saturation rigidity identifies equality
exactly with the presence of every composition of rank $j$. Fullness
of a rank implies fullness of every smaller rank, by downward closure.
The coefficients at degrees zero and one cancel; beyond $N+1$ they
vanish. The words $1^{\len w+1}$ and $(\max w+1)$ give the upper bounds
on $\kappa$, and every nonempty host contains $(1)$.
\end{proof}

\begin{lemma}[Threshold decomposition at a maximum]
Set $\kappa(\varnothing)=1$. For an arbitrary composition
$w=(X,M,Y)$ with $M\ge\max(X,Y)$,
\[
 \kappa(w)=\min\{M+1,\kappa(X)+\kappa(Y)\}.
\]
\end{lemma}
\begin{proof}
The singleton $(M+1)$ is absent. If $u$ and $v$ are missing compositions
of minimum weights for $X$ and $Y$, respectively, then $uv$ is absent
from $XMY$: an embedding of $u$ cannot finish inside $X$, leaving
an embedding of $v$ entirely inside $Y$, a contradiction.

Conversely, take any composition of weight
$k<\min\{M+1,\kappa(X)+\kappa(Y)\}$. If $k<\kappa(X)$ it embeds
in $X$. Otherwise choose its first part for which the cumulative weight
reaches $\kappa(X)$. Its preceding parts have weight below $\kappa(X)$
and embed in $X$; its following parts have weight at most
$k-\kappa(X)<\kappa(Y)$ and embed in $Y$. The chosen part is at most
$k\le M$ and embeds in the central maximum.
\end{proof}

This decomposition uses the same maximum-layer universality mechanism
as Albert, Engen, Pantone and Vatter~\cite{AlbertEtAl}, proof of Theorem~1.
We use it to identify the parameter in the growth
bound; the universality construction is not claimed as new here.
It also gives a linear-time procedure on a maximum Cartesian tree of
the host, and a minimum-weight missing composition as a witness.

\begin{proposition}[Universality threshold for a monotone arm]
If $Y=(t_1\le\cdots\le t_l)$ is nonempty, the largest weight at which
all compositions occur is
\[
 \kappa(Y)-1=\min\left\{l,\ \min_{1\le i\le l}(i-1+t_{l-i+1})\right\}.
\]
\end{proposition}

\begin{proof}
Fix a weight $k$. Necessity of $l\ge k$ follows from $1^k$. For every
$1\le i\le k$, the word $(k-i+1,1^{i-1})$ requires
$t_{l-i+1}\ge k-i+1$. These inequalities are also sufficient: a word
of weight $k$ and length $i$ has every part at most $k-i+1$, and hence
embeds in the last $i$ positions. The inequalities for $i>k$ are
automatic, so the displayed minimum is equivalent to all the necessary
and sufficient conditions.
\end{proof}

The exact support statement identifies the first strict-growth rank
without a monotonicity hypothesis on the host. We next quantify the
deficit at every subsequent rank.
Enumeration and equivalence of single-pattern avoidance classes in
this order are studied by Menon and Singh~\cite{MenonSingh}. The
bounds below apply to every downset inside such an avoidance class,
not only to the full avoidance class.

\begin{proposition}[A quantitative saturation bound]
For every finite composition downset and every $k\ge2$, set
$s=f_k$ and $N=2^{k-1}$. Then
\[
 (k-1)N(2f_k-f_{k+1})\ \ge\ 2s(N-s).
\]
In particular, if $s>0$, then
\[
 \frac{f_{k+1}}{f_k}
 \le 2-\frac{2}{k-1}\left(1-\frac{f_k}{2^{k-1}}\right).
\]
\end{proposition}

\begin{proof}
Put $n=k-1$. Encode feasible words of weight $k$ as a set $S$ of binary
strings of length $n$. A feasible word of weight $k+1$ determines an
edge $x\longrightarrow x_2\cdots x_n b$ between two members of $S$:
the endpoints are obtained by removing the last and the first unit of
the word. The edge determines the word uniquely. Consequently
$2s-f_{k+1}$ is at least the number $e$ of directed edges from $S$ to
its complement in the binary shift graph.

There is a unique directed path of length $n$ from any initial string
$x$ to any terminal string $y$, obtained by appending the bits of $y$.
The $s(N-s)$ paths with $x\in S$ and $y\notin S$ each traverse an edge
from $S$ to its complement. Fix an edge and a position $i\in\{1,\ldots,n\}$
at which that edge is traversed. Without restrictions on the initial
and terminal strings, the number of such length-$n$ paths is
$2^{i-1}2^{n-i}=N/2$. A union bound over the $e$ boundary edges and
the $n$ possible positions gives $s(N-s)\le neN/2$.
Combine this with $e\le2s-f_{k+1}$.
\end{proof}

\begin{theorem}[Sharp uniform deficit from a missing composition]
Suppose a composition of weight $\kappa\ge2$ is absent from a
composition downset. Then, for every $k\ge\kappa-1$,
\[
 2f_k-f_{k+1}\ \ge\ \frac{f_k}{2^{\kappa-1}-1}.
\]
The assertion applies to infinite downsets as well, since each weight
contains only finitely many compositions. The constant is best possible
for every $\kappa$.
\end{theorem}

\begin{proof}
At $k=\kappa-1$, saturation rigidity and the missing word imply
$2f_k-f_{k+1}\ge1$ whenever $f_k>0$, while
$f_k\le2^{\kappa-2}\le2^{\kappa-1}-1$. The zero case is immediate.
It remains to consider $k\ge\kappa$.
Use the preceding binary shift graph on strings of length $n=k-1$,
and put $h=\kappa-1\le n$. Append the $h$ cut bits of a fixed missing
composition $v$. The resulting composition contains $v$: its last
$\kappa$ units have the prescribed internal cuts, and the first part
of $v$ may only be enlarged by preceding units. Hence the terminal
string is outside $S$. Starting from each of the $s=f_k$ members of
$S$ gives $s$ paths of length $h$, each meeting an outgoing boundary
edge. At a fixed position $i$, a fixed edge can be traversed by at
most $2^{i-1}$ of these paths, because the appended block
is fixed. Summing over the $h$ positions and the
$e$ outgoing edges gives $s\le e\sum_{i=1}^h2^{i-1}=e(2^h-1)$. Now use
$e\le2f_k-f_{k+1}$.

For sharpness, take the downset of all compositions with parts at most
$\kappa-1$, and compare weights $k=\kappa$ and $k+1$.
Exactly one composition is excluded at weight $\kappa$, whereas
exactly three are excluded at weight $\kappa+1$, namely
$(\kappa+1)$, $(\kappa,1)$ and $(1,\kappa)$. Thus
$f_\kappa=2^{\kappa-1}-1$ and
$2f_\kappa-f_{\kappa+1}=1$. A finite principal ideal gives the same
example: use the rectangular host $(\kappa-1)^{\kappa+1}$.
\end{proof}

\section{A log-concave cumulative reverse profile}

For a host $Z$ of weight $c$, put
\[
 q_Z(z)=z^cF_Z(1/z),\qquad
 A_Z(k)=\sum_{j=0}^k[z^j]q_Z(z),
\]
with $A_Z(k)=0$ for $k<0$. The cumulative sequence is constant for $k\ge c$.
Write $[t]=1+z+\cdots+z^{t-1}$ for $t\ge1$. For a fixed
legal peak $M\ge\max Z$, set
\[
 D_Z=1+(S_M-1)F_Z,\qquad
 \delta_Z(z)=z^{c+M}D_Z(1/z).
\]
The kernel is nonnegative, since $D_Z=F_{ZM}-F_Z$.
Parenthesized integer arguments of $D_Z$ and $\delta_Z$ below denote
coefficients, not polynomial evaluations. We first prove the two
cumulative properties when $Z$ is increasing.

\begin{theorem}
For every increasing arm $Z$, the sequence $A_Z$ is log-concave:
\[
 A_Z(k)^2\ge A_Z(k-1)A_Z(k+1)\qquad(k\in\mathbb Z).
\]
In particular, for a legal peak $M$, the reversed kernel $\delta_Z$ is
log-concave at each index $0\le k\le M-2$.
\end{theorem}

\begin{proof}
The empty arm gives $A_\varnothing(k)=\mathbf1_{k\ge0}$, which is
log-concave with interval support. Suppose the result holds for an arm of
weight $c$ and append a global maximum $t\ge1$. The reversed recurrence is
\[
 q_{Zt}=[t]q_Z+z^{c+t}.
\]
Consequently
\[
 A_{Zt}(k)=W(k)+\mathbf1_{k\ge c+t},\qquad
 W(k)=\sum_{i=0}^{t-1}A_Z(k-i).
\]
The convolution of two nonnegative log-concave sequences with interval
support is log-concave; it applies here to $A_Z$ and $[t]$. Although $A_Z$
has infinite constant support to the right, this application is local:
truncate it strictly after all indices occurring in any of the three
coefficients being compared. Such a terminal truncation preserves
log-concavity and interval support.

Set $C=tA_Z(c)$. Since $[z^c]q_Z=[z^0]F_Z=1$, the last two values before
the added step are
\[
 W(c+t-2)=C-1,\qquad W(c+t-1)=C,
\]
and $W(k)=C$ thereafter. These formulas also cover $c=0$ and $t=1$, with
zero extension. All log-concavity inequalities before index $c+t-1$
remain unchanged. The boundary values are $C-1,C,C+1,C+1,\ldots$.
The only two remaining boundary inequalities are
\[
 C^2-(C-1)(C+1)=1>0,\qquad
 (C+1)^2-C(C+1)=C+1>0.
\]
Afterward the sequence is constant. Induction proves the first assertion.

The kernel identity is
\[
 \delta_Z=([M]-z^M)q_Z+z^{c+M}.
\]
Thus $[z^k]\delta_Z=A_Z(k)$ for $0\le k<M$. All three coefficients in a
log-concavity comparison agree when $0\le k\le M-2$, including the zero
coefficient at index $-1$. This proves the second assertion. No extension
past that interval is asserted.
\end{proof}

\begin{theorem}[Cumulative reverse kernel]
Fix $M\ge2$. For every increasing arm $Y$ with parts strictly below $M$,
the sequence
\[
 B_Y(k)=\sum_{j=0}^k[z^j]\delta_Y(z),\qquad B_Y(k)=0\ (k<0),
\]
is log-concave on all integer indices.
\end{theorem}

\begin{proof}
The empty arm gives $\delta_\varnothing=[M]$, so the assertion is clear.
Suppose $Y=Zt$, put $c=\wt Z$, and let $b=c+t$. The reversed recurrence
gives
\[
 B_Y(k)=W(k)+\sum_{j=0}^k[z^j]\bigl(z^b[M-t]\bigr),\qquad
 W=[t]*B_Z.
\]
Induction and the log-concave convolution lemma make $W$ log-concave.
For a comparison centered at $k\le b-2$, the correction has not started,
and all three coefficients agree with $W$.

For $k\ge b$, the kernel's known low-rank growth says
$\delta_Y(k+1)\le\delta_Y(k)$: after reversing, these are forward ranks
at most $M$. Consequently $B_Y$ is nonnegative, nondecreasing and
concave at such a comparison, which implies log-concavity there.
Indeed, if consecutive increments are $s\ge t\ge0$ and the first value
is $A\ge0$, the log-concavity determinant is $s^2+A(s-t)\ge0$.

Only $k=b-1$ remains. It is covered by the same argument if
$\delta_Y(b)\le\delta_Y(b-1)$. Otherwise the reversed recurrence gives
\[
 0<\delta_Y(b)-\delta_Y(b-1)
   =1+D_Z(M-t)-D_Z(M).
\]
The two forward coefficients are integers and $D_Z$ is nondecreasing
through $M$, so they must be equal. In the kernel proof, the identity
\[
 (1-z)D_Z=E_Z-z^{M+1}F_Z
\]
and the positive interval support $\supp E_Z=[q,c+1]$, with $q\le t+1$,
then force $c+1\le M-t$: otherwise the interval $[M-t+1,M]$ meets the
positive support and $D_Z(M)>D_Z(M-t)$. Hence $b<M$.

For every $0\le j\le b$, therefore, $\delta_Y(j)=A_Y(j)$ by the preceding
theorem. The three values of $B_Y$ at the remaining comparison are
cumulative sums of the log-concave sequence $A_Y$. Cumulative summation
preserves log-concavity, by convolution with the sequence
$\mathbf1_{k\ge0}$ (or a sufficiently long finite window). This proves
the final comparison and completes the induction.
\end{proof}


\section{Tail counts and composition at a common maximum}\label{sec:cumulativecores}

We distinguish log-concavity of a coefficient sequence from log-concavity
of its tail sums. For a nonnegative polynomial $P$ of degree $d$, write
\[
 \mathcal T_P(k)=\sum_{j=k}^{d}P_j\qquad(k\in\mathbb Z),
\]
where coefficients outside $[0,d]$ vanish. Thus $\mathcal T_P$ is constant
for $k\le0$ and zero for $k>d$.

\begin{lemma}[Cumulative convolution closure]
Let $p,q$ be nonzero nonnegative finitely supported sequences. If their
cumulative sums are log-concave on all integer indices, the cumulative
sums of $p*q$ are log-concave as well.
\end{lemma}
\begin{proof}
Normalize $p$ and $q$ by their positive total masses and translate their
supports into the nonnegative integers. The hypothesis says precisely
that both distributions are discrete decreasing reversed hazard
distributions. Apply Nanda and Sengupta~\cite{NandaSengupta},
Theorem 4.2(d). Undo the normalization and the translations.
\end{proof}

\begin{lemma}[Closure at a common maximum]\label{lem:cumulativegraft}
Fix $M\ge2$ and arbitrary hosts $X,Y$ with parts at most $M$.
Suppose the cumulative sums of each of $q_X,q_Y,\delta_X,\delta_Y$
are log-concave, where the two kernels use peak $M$.
Then both the reversed rank polynomial and the reversed peak-$M$
kernel of $XMY$ have log-concave cumulative sums.
\end{lemma}
\begin{proof}
The maximum identity gives
\[
 D_{XMY}=D_XD_Y,\qquad
 q_{XMY}=q_X\delta_Y+z^{a+M}q_Y,
 \qquad a=\wt X,\quad b=\wt Y.
\]
The first identity follows by expanding
$1+(S_M-1)(F_X+F_Y+(S_M-1)F_XF_Y)$.
After reversing, it gives $\delta_{XMY}=\delta_X\delta_Y$;
cumulative convolution closure proves the kernel assertion.

For the rank polynomial, interchange the arms if necessary so that
$a\ge b$. This leaves the polynomial unchanged. The cumulative sums of
$q_X\delta_Y$ are log-concave. At centers $r\le a+M-2$, the
added term has not entered any of the three cumulative values.
At every remaining center, the consecutive increments of the full
cumulative sequence satisfy
\[
 (q_{XMY})_r\ge(q_{XMY})_{r+1}.
\]
Indeed, their forward ranks are $a+b+M-r\le b+1$ and the preceding
rank, while $b+1\le\lceil(a+b+M)/2\rceil$. Apply
Lemma~\ref{lem:half}; ranks outside the support cause no difficulty.
Nonnegative decreasing increments imply concavity and hence
log-concavity of the cumulative sequence. This includes the joining
center $a+M-1$.
\end{proof}

\begin{theorem}[Valleys separated by a common maximum]
\label{thm:separated}
Let
\[
 w=Z_0\,M\,Z_1\,M\cdots M\,Z_s,
\]
where each $Z_i$ weakly decreases and then weakly increases, possibly
with one or both segments empty, and every part of $Z_i$ is strictly
below $M$.
There is no restriction on $M$, $s$, or the block lengths.
Both $F_w$ and $1+(S_M-1)F_w$ have log-concave tail counts.
In particular, the number of descents need not be bounded.
\end{theorem}
\begin{proof}
For $M=1$, all blocks are empty, and the word consists of ones;
both conclusions follow directly from the corresponding intervals.
Assume $M\ge2$. Every nonempty valley block has a global maximum at
an endpoint; removing that maximum leaves a valley block. Word reversal
preserves the rank polynomial, so removing either endpoint maximum $t$
uses the identity $F=1+S_tF_{\mathrm{rest}}$. Induction shows that a
valley block has the same rank polynomial as its increasing rearrangement.
Thus each block has the two cumulative properties proved for monotone
arms in the preceding section. Repeated application of
Lemma~\ref{lem:cumulativegraft} proves both conclusions. Equivalently,
the kernel factorization is
\[
 D_w=\prod_{i=0}^s D_{Z_i}.
\]
This factorization concerns the kernels at the same peak $M$.
\end{proof}

For example, $(1,7,10,8,3,10,2,9,10,7,4,1)$ has five descents and
satisfies the theorem with $M=10$. The theorem remains a statement about
cumulative counts; it makes no rank-unimodality assertion for this
larger class. Appending or prepending successive global maxima also
preserves the rank-tail conclusion, by the reverse append argument.

\begin{theorem}[Log-concavity of one-descent tail counts]
Let $w$ be any composition with at most one adjacent descent, with no
bound on its length or part sizes. Then
\[
 \mathcal T_{F_w}(k)^2\ge
 \mathcal T_{F_w}(k-1)\mathcal T_{F_w}(k+1)
 \qquad(k\in\mathbb Z).
\]
The same conclusion holds for every bitonic composition. For increasing
strict arms $X,Y$, it also holds for the relative polynomial $F_XD_Y$.
\end{theorem}

\begin{proof}
Consider first a strict core $(X,M,Y)$. The proved cumulative reverse-arm and reverse-kernel
theorems, together with cumulative convolution closure, show that the
cumulative sums of $q_X\delta_Y$ are log-concave. This proves the
relative-polynomial assertion, by reversing the index.
The core itself is covered by Theorem~\ref{thm:separated} with two
blocks.

Monotone hosts are already covered by the reverse-arm theorem.
Appending a global maximum changes $F$ to $1+S_tF$ and preserves
log-concavity of the cumulative reverse profile by the same window and
terminal-step argument used for monotone arms. That argument requires
only a nonnegative polynomial with constant coefficient one, a
log-concave cumulative reverse profile, and the maximum append identity.
Strip successive global maxima from the terminal increasing segment of
a one-descent word until the strict core is reached, and then restore
them. Inserting another copy next to the central peak also changes
$F$ to $1+S_MF$: its peak-$M$ kernel is multiplied by $S_M$, by the
maximum identity. Thus repeated copies of the peak are restored by
the same cumulative argument.
Finally, reversing the decreasing arm of a bitonic word does not change
its arm polynomial; the maximum decomposition gives the identical
core polynomial. This proves all assertions.
\end{proof}

The theorem asserts log-concavity of tail \emph{counts}, not tail
log-concavity of the individual rank coefficients. The latter is a
different, stronger shape condition. Log-concavity of tail sums alone
does not imply rank-unimodality.

\paragraph{External input and verification boundary.}
The cumulative convolution lemma is an external theorem. Ordinary
convolution closure for log-concave coefficient sequences does not
justify this step. The supplied Lean files verify the reverse append
and escape counting models; they do not formalize this external
probability theorem or the complete core theorem.

\section{Arbitrary hosts with a dominant maximum}\label{sec:dominant}

The maximum decomposition also gives a length-unrestricted family with
no restriction on the number of descents.

\begin{lemma}[Two words at each internal rank]\label{lem:twointernal}
If a host $w$ has at least two parts and has a part greater than one,
then $[z^k]F_w\ge2$ for $2\le k\le\wt w-1$.
\end{lemma}
\begin{proof}
Put $l=\len(w)$ and $n=\wt w$. If $k\le l$, use $1^k$ and a word
consisting of one two and $k-2$ ones, embedded using a non-one part.
If $k>l$, choose positive integers $c_i\le w_i$ with sum $k$.
If there are two such vectors, they give two distinct words. Otherwise,
choose a slack coordinate $j$, so $c_j<w_j$; this exists because $k<n$.
Every other $c_i$ equals one, since transferring a unit from $c_i>1$
to $c_j$ would give a second vector. Delete one of these other ones and
increase $c_j$ by one. This gives a second word of weight $k$, now of
length $l-1$.
\end{proof}

\begin{lemma}[A comparison at distance two]\label{lem:codimtwo}
For an arbitrary host of weight $n\ge5$, write $F_w=\sum p_i z^i$.
Then $p_{n-2}\ge p_3\ge p_2$.
\end{lemma}
\begin{proof}
The second inequality is Lemma~\ref{lem:half}. If $p_3\le2$, the first
follows from Lemma~\ref{lem:twointernal}, or directly when the host is
a single part or consists only of ones. Assume $p_3\ge3$. Write $e$
for the number of non-one parts and $h$ for the number of parts at
least three. Distinct words obtained by subtracting two units without
deleting a part number exactly $\binom e2+h$. These have the original
length; deleting one or two parts gives disjoint additional counts.

If $e\ge3$, this count is at least three. When $p_3=4$, the word $(3)$
occurs, so $h\ge1$, and the count is at least four.

Suppose $e=2$ and $h\ge1$. If both parts are at least three, the
same-length count is three. If a one is present, decrementing a
non-one part and deleting a one supplies another word. If no one is
present, $p_3\le3$, which is enough. If just one non-one part is at
least three, the same-length count is two. With no ones, deleting the
part two supplies a third word and $p_3\le3$. With ones present,
decrementing either non-one part and deleting a one gives two further
words: one retains two non-one parts and the other only one.

If $e=2$ and $h=0$, then $p_3\le3$ and at least one one is present.
The same-length count is one. If there are at least two runs of ones,
delete a one from each of two runs, and separately delete a part two;
these add two words of different lengths. If there is one run between
the two twos, deleting either two gives two distinct further words.
If the run is at an endpoint, reverse the host if needed and write it
as $1^r22$. Deleting a two gives $1^r2$, while deleting a one and
decrementing the final two gives $1^{r-1}21$. Thus the count is again
at least three.

It remains to consider $e=1$, with non-one part $t$.
If $t\ge3$ and the ones are all on one side, $p_3\ge3$ requires at
least two ones. Subtracting two from $t$, decrementing $t$ and deleting
a one, and deleting two ones give three different lengths; here
$p_3\le3$. If ones occur on both sides, subtracting two from $t$ and
deleting one from each side supply two words of different lengths.
Decrementing $t$ and deleting a one on either side supplies two more,
distinguished by the number of ones before the surviving non-one part.
Finally, if $t=2$ and $p_3\ge3$, ones occur on both sides. Since
$n\ge5$, one of the two runs has length at least two. Deleting two
ones from that run, or one from each run, gives two distinct words;
deleting the two gives a third word of a different length. Here
$p_3=3$. These cases exhaust the possibilities.
\end{proof}

\begin{lemma}[Central product inequality]\label{lem:centralproduct}
Let $A=F_X$, $B=F_Y$, $a=\wt X$, $b=\wt Y$, $D=a+b\ge6$, and
$C=AB=\sum c_i z^i$. Then
\[
 \Del A(D-1)+\Del B(D-1)+2c_{D-2}-c_{D-1}-c_3\ge0.
\]
Moreover $c_3\ge c_2$.
\end{lemma}
\begin{proof}
Interchange the arms so that $a\ge b$.
For a nonempty arm $Z$ of weight $n$, put
$q=p_{n-1}$, $r=p_{n-2}$ and $L=2r-p_2-p_3$.
Then $q\ge1$. Lemma~\ref{lem:codimtwo} gives $L\ge0$ for $n\ge5$.
For $n=4$ one has $L=r-q\ge1-q$, and for $n=3$ one has
$L=1-q$. Thus $L\ge1-q$ whenever $n\ge3$.

If $b=0$, the displayed expression is $r-p_3\ge0$ for the other arm.
If $b=1$, it is $L\ge0$, since $a\ge5$. If $b=2$, it is $L+q-1$;
here $a\ge4$, so this is nonnegative. If $b\ge3$, write $s,t$ for
the quantities $q,r$ of the second arm, and $L_B=2t-B_2-B_3$.
Neither arm contributes a boundary difference at rank $D-1$, and the
expression is
\[
 L+L_B+2qs-q-s
 \ge 2(q-1)(s-1)\ge0.
\]
These identities follow by expanding the top two and bottom three
coefficients of $AB$; all nonempty arm polynomials have constant and
linear coefficient one, and leading coefficient one.

For the last assertion, if one arm is empty use
Lemma~\ref{lem:half} for the other, whose weight is at least six.
Otherwise $c_3-c_2=A_3+B_3-1\ge0$, since at least one arm has
weight at least three.
\end{proof}

\begin{theorem}[Dominant maximum]\label{thm:dominant}
Let $w$ be an arbitrary nonempty composition of weight $N$ and maximum
part $M$. If $N\le2M+5$, then $F_w$ is unimodal.
\end{theorem}
\begin{proof}
The case $M=1$ is an interval polynomial. Assume $M\ge2$, split at
any maximum as $w=XMY$, and put $D=\wt X+\wt Y=N-M$.
The cases $D=0$ and a single empty arm cause no change to the identities
below. With $A=F_X$, $B=F_Y$, $C=AB$, the maximum decomposition gives
\begin{equation}\label{eq:dominantdifference}
 \Del F_w(n)=\Del A(n)+\Del B(n)
       +2c_{n-1}-c_n-c_{n-M-1}.
\end{equation}
Consequently all differences at $n\ge D+2$ are nonpositive.
If $M<D$, the difference at $D+1$ is also nonpositive:
when both arms are nonempty it equals $2-c_{D-M}\le0$, because
$c_j\ge2$ for $1\le j<D$; with an empty arm it equals
$1-c_{D-M}\le0$.

If $M\ge D-1$, Lemma~\ref{lem:half} gives increase through rank $D$,
so at most the single difference at $D+1$ remains undecided before
the permanently decreasing part. If $M=D-2$ or $M=D-3$, it gives
increase through rank $D-1$; the difference at $D$ is the only
undecided one. These cases are unimodal.

The two remaining cases are $M=D-4$ and $M=D-5$. Here $D\ge6$,
and first-half increase reaches $D-2$. At $n=D-1$, the subtracted
coefficient in~\eqref{eq:dominantdifference} is respectively $c_2$
or $c_3$. Lemma~\ref{lem:centralproduct} makes this difference
nonnegative in both cases. Again only the difference at $D$ remains
undecided, followed by nonpositive differences at every larger rank.
This proves unimodality.
\end{proof}

This theorem allows arbitrary lengths and descents. Its condition is
a sufficient one; no optimality of the constant five is asserted.


\section{Minimal counterexamples and outer layers}\label{sec:peeling}
The following reduction reconstructs the earlier peeling argument, with all hypotheses needed here included.

\begin{lemma}[Balanced minimal counterexamples]\label{lem:balanced}
If, for fixed $M$, some $Q_{X,Y}$ is not unimodal, choose one with the fewest total arm parts. It has nonempty arms and $|a-b|\le M-1$. Within this band, $Q_{X,Y}$ is unimodal if and only if $H$ is unimodal, so its arms may be exchanged. More generally, this equivalence holds whenever $b\le a+M-1$. Both arm polynomials in the minimal counterexample may be assumed nonrectangular, and we may orient it with $a\ge b$.
\end{lemma}
\begin{proof}
Empty arms give either $D_Y$ or $S_MF_X$, both unimodal. Every shorter $Q$ is unimodal by minimality. If $a\ge b+M$, the guaranteed growth endpoint in~\eqref{eq:qgrowth} is at least $b+M+1$. At that index and beyond, the correction $D_Y$ in the first recurrence~\eqref{eq:qappend} has vanished; the sequence joins the unimodal $S_uQ_{Z,Y}$. This is impossible for a counterexample. If $b\ge a+M$, the second recurrence works identically, its correction having degree $a+M$.

When $b\le a+M-1$, both $H$ and $Q$ increase through rank $b+1$, by Lemmas~\ref{lem:half} and~\ref{lem:kernel}; their differences agree from $b+2$ onward. They are therefore unimodal simultaneously. Both directions satisfy this condition in the balanced band. Finally, an interval $F_X$ times $D_Y$ is unimodal, and the case of interval $F_Y$ follows by exchanging the arms.
\end{proof}

\begin{lemma}[A splice criterion]\label{lem:splice}
Write $X=Zu$, $R=b+M$, $V=Q_{Z,Y}$, and $P=S_uV$. Suppose $V$ is unimodal and $Q$ is nondecreasing through $R$. If $Q$ is not unimodal, necessarily
\begin{equation}\label{eq:plateau}
 V_{R-u}=V_{R-u+1}=\cdots=V_R,
 \qquad P\text{ increases strictly at a later rank}.
\end{equation}
Thus either strict inequality $V_j>V_{R-u}$ for some $R-u<j\le R$, or two unequal successive differences inside this interval, rules out this obstruction.
\end{lemma}
\begin{proof}
At the splice $\Del Q(R+1)=V_R-V_{R-u}-1=\Del P(R+1)-1$; later differences equal those of $P$. Since $P$ is unimodal and differences are integers, failure requires $\Del P(R+1)=0$ with later increase. The equal endpoints of a unimodal $V$ must then enclose a constant interval. If an internal value were larger, $V$ would already have passed a peak, and each subsequent sliding-window difference would be nonpositive. This proves~\eqref{eq:plateau}.
\end{proof}

\begin{lemma}[Strict reflected gap]\label{lem:gap}
Let $\delta\ge2$, $b\ge\delta$, $\delta\le u<M$, and $\wt Z=R-\delta-u\ge0$. Then
\[
 (F_ZD_Y)_{R-\delta+1}>(F_ZD_Y)_{R-u}.
\]
\end{lemma}
\begin{proof}
All products of layers have centers at least $(\wt Z+R+1)/2$, the midpoint of the two indices being compared. Their contributions are nonnegative. Lemma~\ref{lem:budget} at $t=\delta-1$ gives $d_{R-\delta+1}>d_\delta$. Thus the layer of $D_Y$ at height $d_\delta+1$ has left endpoint $l\ge\delta+1$ and right endpoint $r\ge R-\delta+1$. Multiply this layer by $\Iv0{\wt Z}$. At the two indices, its coefficients are
\[
 \max\{0,\min(\wt Z,n-l)+1\},\qquad n=R-u,\ R-\delta+1.
\]
The latter is strictly larger: the earlier index is at most $\wt Z+l-1$, the later is at least $l$, and the indices differ by $u-\delta+1\ge1$.
\end{proof}

\begin{proposition}[The outer three layers]\label{prop:outer}
Every minimal counterexample as in Lemma~\ref{lem:balanced}, for $M\ge3$, satisfies $a-b\le M-4$. For $M=2$ there are no counterexamples.
\end{proposition}
\begin{proof}
For $M=2$ all arm polynomials are intervals. Otherwise retain the nonrectangular, heavier-left orientation, so $b\ge3$ and $u=\max X\ge2$.

If $a=R-1$, growth is known through $R$. In $S_uF_ZD_Y$, all layer-product centers are at least $R+1/2$, so every contribution to $\Del P(R+1)$ is nonnegative. Let $[q,R]$ be the lowest layer of $D_Y$, where $q\ge2$. Its product with the lowest layer of $F_Z$ contributes
\[
 \coef{R-q}{[u][R-u]}\ge1,
 \qquad [t]=1+\cdots+z^{t-1}.
\]
Here $0\le R-q\le R-2$, exactly the support of this product. Hence $\Del Q(R+1)\ge0$, and the unimodal tail joins.

If $a=R-2$, growth again reaches $R$. Apply Lemma~\ref{lem:gap} with $\delta=2$ to rule out~\eqref{eq:plateau}.

Suppose $a=R-3$. Growth initially reaches $R-1$; we first gain one further rank. Each layer-product in $P=S_uF_ZD_Y$ has center at least $R-1/2$. The $d_{R-1}-1$ layers of $D_Y$ of height at least two that reach rank $R-1$ have right endpoint $R-1$ and left endpoint between $3$ and $R-1$, since $d_2\le1$ and $d_R=1$. Each, multiplied by the lowest layer of $S_uF_Z$, contributes at least one to $\Del P(R)$. It follows that
\[
 \Del Q(R)=1-d_{R-1}+\Del P(R)\ge0.
\]
If $u\ge3$, Lemma~\ref{lem:gap} with $\delta=3$ finishes the splice.

It remains that $u=2$, with $\alpha=\wt Z=b+M-5$. We show that $\Del V(R-1)>0$ when $b\ge4$. Put $A=F_Z=\sum a_jz^j$, $B=F_Y$, and $C=AB$. Since $\deg A=R-5$,
\[
 \Del V(R-1)=2c_{R-2}-c_{R-1}-c_{b-2}.
\]
Decompose $B$ into layers. For its lowest layer $[0,b]$, put $L=M-3$. Telescoping gives the contribution
\[
 \sum_{j=\alpha-L+1}^{\alpha}a_j-\sum_{j=0}^{L}a_j+a_{L+1}\ge0.
\]
Indeed $\sum_{j=1}^La_j\le\sum_{j=\alpha+1-L}^{\alpha}a_j$ by strengthened reflection, and $a_{L+1}\ge a_0=1$. For a higher layer $[l,r]$, put $U=b-l$ and $v=b-r$. Its contribution is exactly
\[
 d^Z_{U+M-1}-d^Z_{v+M-2},\qquad D_Z=\sum d^Z_jz^j.
\]
The indices are ordered and their sum is at most $b+2M-4=\alpha+M+1$, so the contribution is nonnegative. The second layer $[2,b-1]$ contributes $d^Z_{\alpha+2}-d^Z_{M-1}\ge a_{M-1}\ge1$ by Lemma~\ref{lem:budget}. This excludes a plateau.

If $b=3$, then $\deg C=R-2$ and $c_1=2$. The same identity gives $\Del V(R-1)=0$, while $\Del V(R)=-c_2<0$. Again there is no plateau. All outer cases are excluded.
\end{proof}



\section{Finite profiles for unbounded arms}\label{sec:profiles}
For a polynomial $F_Z$ of degree $c$, write $F_Z\rev=z^cF_Z(1/z)$. The results in this section permit any fixed part bound. No length bound is assumed.

\begin{lemma}[Deletion of an initial part]\label{lem:prefixreverse}
If $t\le\min U$, then $F_{tU}\rev-F_U\rev$ is coefficientwise nonnegative.
\end{lemma}
\begin{proof}
Let $u=\wt U$, $\RR_U=\prod_{v\in U}S_v$, $T_U=\prod_{v\in U}[v]$, and
\[
 K_U=(1-z)F_U+z\RR_U.
\]
The polynomial $K_U$ counts the empty word and the smaller words whose first letter is at least two. To verify this interpretation, both sides satisfy $K_{Ut}=1-z+S_tK_U$ with initial value one. Thus $K_U\ge0$ and its degree is at most $u$. Expanding the monotone-arm recurrence at the initial part gives $F_{tU}=F_U+S_t\RR_U$. On reversing, this yields
\[
 F_{tU}\rev=z^tF_U\rev+[t]T_U.
\]
The definition of $K_U$ also gives $zK_U\rev=T_U-(1-z)F_U\rev$. Consequently
\begin{equation}\label{eq:prefixpositive}
 F_{tU}\rev-F_U\rev=z[t]K_U\rev\ge0.
\end{equation}
The argument permits $t=1$ and empty $U$.
\end{proof}

\begin{lemma}[Capping multiplicities]\label{lem:capping}
Replace each multiplicity in an increasing arm $Z$ by its minimum with an integer $L\ge1$, obtaining $\bar Z$.
\begin{enumerate}
\item The coefficients of $F_Z$ and $F_{\bar Z}$ agree through degree $L$.
\item For $k<L$, $(F_Z\rev)_k\ge(F_{\bar Z}\rev)_k$.
\item For a fixed legal peak $M$, put $\delta_Z=z^{\wt Z+M}D_Z(1/z)$. For $k<L$, $(\delta_Z)_k\ge(\delta_{\bar Z})_k$.
\end{enumerate}
\end{lemma}
\begin{proof}
A smaller word of weight at most $L$ has at most $L$ letters. In an increasing host, any embedding of such a word can be moved to the last $L$ positions, or all positions if there are fewer. Capping retains those positions, so it changes no such word.

The exact reversed recurrences are
\begin{equation}\label{eq:reverseappend}
 F_{Zt}\rev=[t]F_Z\rev+z^{\wt Z+t},\qquad
 \delta_{Zt}=[t]\delta_Z+z^{\wt Z+t}[M-t].
\end{equation}
Reinsert removed copies after the $L$ retained copies in their equal-letter block. At the first reinsertion the prefix weight is at least $L$. Thus every new correction monomial has degree greater than $L$, including those at later append steps. Modulo $z^L$, reinsertion only multiplies the already nonnegative reversed polynomial by $[t]$, whose constant coefficient is one. Later append operations also preserve coefficientwise inequalities. Repeating this for each clipped block proves both comparisons.
\end{proof}

\begin{theorem}[A general high-layer criterion]\label{thm:generalhigh}
Let $m,r,N$ be integers with $m,r\ge2$ and $N\ge r$. Put
\[
 J=(r-1)m,\qquad K=\max(0,J-r).
\]
Assume
\begin{equation}\label{eq:binomialcondition}
 \binom Nk\ge2^{r+k-1}\quad(1\le k\le K).
\end{equation}
If an increasing arm $Y$ has parts at most $m$, at least $N$ non-one parts, weight $b$, and coefficients $p_j$, every layer of height greater than $p_r$ has endpoint sum at least $b+r$. A convenient sufficient condition for~\eqref{eq:binomialcondition} is
\begin{equation}\label{eq:generalcutoff}
 N\ge\max\{r,2^r,3K-1\}.
\end{equation}
\end{theorem}
\begin{proof}
Since $b\ge2N\ge2r$, the coefficients through rank $r$ are nondecreasing. Peel the last $r$ parts. The append recurrence expresses $F_Y=B+G$, where $\deg B\le(r-1)m=J$, $B\ge0$, and $G$ is the product of the remaining arm polynomial with the $r$ corresponding $S_t$ factors. Its interval-product summands have centers at least $(b+r)/2$. Thus reflection about $(b+r)/2$ holds at every left index $j>J$.

For $r<j\le J$, put $k=j-r$. Expansion of the reversed recurrence~\eqref{eq:reverseappend} contains the product of all $[t]$ factors as a nonnegative summand. Every non-one factor contains $1+z$, so
\[
 p_{b+r-j}=(F_Y\rev)_k\ge\binom Nk\ge2^{j-1}\ge p_j.
\]
The last inequality counts all compositions of weight $j$. Indices $j\le r$ cannot start a layer of height greater than $p_r$. The same left-endpoint contradiction as in~\eqref{eq:comparison} proves the layer assertion.

For the sufficient condition, $\binom N1=N\ge2^r$. For $2\le k\le K$, the ratio of successive binomial coefficients is $(N-k+1)/k\ge2$, because $N\ge3K-1$. Induction gives~\eqref{eq:binomialcondition}. If $K=0$, there are no binomial inequalities to check.
\end{proof}




\section{A quadratic obstruction bound for arbitrary peaks}\label{sec:uniform}
The profile method has an analytic consequence independent of the finite certificates. The key is to move the height threshold farther to the right, then pay for every potentially negative low layer with one selected positive contribution.

\begin{lemma}[High layers beyond the recurrence error]\label{lem:errorhigh}
Let $m,r\ge1$ be integers. Let $Y$ be an increasing arm of weight $b$, with parts at most $m$, and suppose it has at least $r$ parts. Take an integer $J\ge(r-1)m$ and assume $2J\le b$. Every height layer of $F_Y$ above height $(F_Y)_J$ has endpoint sum at least $b+r$.
\end{lemma}
\begin{proof}
Peeling the last $r$ parts gives $F_Y=B+G$, where $B\ge0$ has degree at most $(r-1)m\le J$, and $G$ is the remaining arm polynomial multiplied by the last $r$ factors $S_t$. The interval-product summands of $G$ have centers at least $(b+r)/2$. Thus $p_j\le p_{b+r-j}$ at every left index $j>J$. First-half increase gives $p_j\le p_J$ for $j\le J$. A layer above $p_J$ with endpoint sum less than $b+r$ would start at $l>J$ and omit its reflected index $b+r-l$, contrary to this comparison.
\end{proof}

\begin{theorem}[Quadratic obstruction bound]\label{thm:uniformcentral}
For $M\ge4$, set
\[
 J_M=\max\{M+1,(M-3)(M-1)\},\qquad
 N_M=8(J_M+M-3)=\max\{48,8M(M-3)\}.
\]
Let $X,Y$ be increasing strict arms of weights $a=b+d$, where $0\le d\le M-4$. If each arm has at least $N_M$ non-one parts, then
\[
 \Del H(n)>0\qquad
 \left(\left\lceil\frac{a+b+M}{2}\right\rceil<n\le b+M+1\right).
\]
Consequently, for any fixed peak $M$, a minimal counterexample to unimodality of $Q$ has an arm with fewer than $\max\{48,8M(M-3)\}$ non-one parts.
\end{theorem}
\begin{proof}
Write $J=J_M$, $K=J+M-3$, $N=N_M=8K$, and $n=b+M+h$ in the stated interval. Then
\[
 h\le1,\quad 2h\ge d-M+2,\quad n-1>a.
\]
Hence the difference of $H$ equals that of $F_YD_X$ here. By Lemma~\ref{lem:strongkernel}, every layer of $D_X$ has endpoint sum at least $a+M+3$. Write $D_X=\sum d_i z^i$ and $\delta=D_X\rev$.

Apply Lemma~\ref{lem:errorhigh} with $r=M-2$ and $m=M-1$; its assumptions hold because $b\ge2N\ge2J$. Every layer $[l,r']$ above height $p_J=(F_Y)_J$ has endpoint sum at least $b+M-2$. Its contribution is $d_{n-l}-d_{n-r'-1}$, whose ordered indices have sum at most
\[
 2n-(b+M-2)-1=b+M+2h+1\le a+M+3.
\]
These high-layer contributions are all nonnegative.

There are exactly $p_J$ remaining layers. Ordinary strengthened reflection gives $l+r'\ge b+1$ for each nonlowest layer, so its negative-term index satisfies
\[
 n-r'-1\le M+h+l-2\le M+J-1.
\]
For the lowest layer $[0,b]$, that index is $M+h-1\le M$, also within the bound. For every integer $i\le M+J-1$, nonnegativity and the universal composition count give
\[
 0\le d_i\le 2^{M+J-2}.
\]
Indeed $d_i\le\sum_{t=1}^M (F_X)_{i-t}\le2^{i-1}$ for $i\ge1$, and $d_i=0$ for $i\le0$.

There is a layer beginning at $J$. To see this, every smaller word of weight $J-1$ embeds in the last $J-1$ positions of $Y$; prefixing a one gives an injection into words of weight $J$. The additional word $2\,1^{J-2}$ embeds because $Y$ has at least $J$ non-one parts and is not in that image. Thus $p_J>p_{J-1}$. The positive term of a selected layer beginning at $J$ is
\[
 d_{n-J}=\delta_k,\qquad k=J+d-h.
\]
The inequalities on $h,d$ give
\[
 J-1\le k\le J+M-3=K.
\]
The reversed kernel recurrence contains the product of the $[t]$ factors as a nonnegative summand. Therefore $\delta_k\ge\binom Nk$. Since $N\ge8k$, selecting one element from each of $k$ disjoint eight-element blocks proves $\binom Nk\ge8^k$. Retain this one positive contribution and bound all the negative contributions of the low layers. As $p_J\le2^{J-1}$, this yields
\[
 \Del H(n)\ge8^{J-1}-2^{J-1}2^{M+J-2}
       =2^{3J-3}-2^{M+2J-3}>0,
\]
where strictness follows from $J\ge M+1$.

For the last assertion, Proposition~\ref{prop:outer} puts a minimal counterexample in the stated weight band. The initial kernel reflection makes $Q$ nondecreasing through $\lceil(a+b+M)/2\rceil$. If both non-one counts were at least $N_M$, the strict interval just proved would extend this through $b+M+1$. The shorter-arm unimodal tail would then join, a contradiction.
\end{proof}

The bound concerns one arm of a minimal obstruction, not its total weight or the size of the other arm. It does not prove the unrestricted conjecture. Unlike the general binomial criterion applied at height $p_{M-2}$, the present argument needs only quadratically many non-one parts: moving the threshold to $p_{J_M}$ makes all higher layers automatically nonnegative, and one low layer pays the entire remaining cost.

\section{A clipped central budget}\label{sec:central}
We first establish the high-layer reflection used uniformly at peaks four through seven. The extra unit of reflection, compared with the ordinary arm bound, is essential to the central comparison.

\begin{theorem}\label{thm:highlayers}
Let $Y$ have parts at most six and at least nine non-one parts, and write $F_Y=\sum p_jz^j$, $b=\wt Y$. Every height layer above height $p_7$ has endpoint sum at least $b+5$.
\end{theorem}
\begin{proof}
Let $N$ be the number of non-one parts. If $N\ge32$, Theorem~\ref{thm:generalhigh} applies with $m=6,r=5$: the nineteen inequalities $\binom{32}{k}\ge2^{k+4}$, $1\le k\le19$, all hold. It gives the stronger assertion for heights above $p_5$.

For $9\le N\le31$, peel the last five parts. The append recurrence gives $F_Y=B+G$, where $B$ is the sum of the empty product and the last one, two, three and four $S_t$ factors, and $G$ is the product of the last five factors with the remaining arm polynomial. Thus $\deg B\le24$, and every interval-product summand of $G$ has center at least $(b+5)/2$. For left indices $j>24$, reflection follows because $B_j=0$ and $B\ge0$. For $j\le7$ the inequality $p_j\le p_7$ follows from first-half increase. It remains to prove
\begin{equation}\label{eq:mirrorcertificate}
 8\le j\le24,\quad 2j\le b+5,\quad p_j>p_7
 \quad\Longrightarrow\quad (F_Y\rev)_{j-5}\ge p_j.
\end{equation}
Cap all multiplicities at 24. Low coefficients through degree 24 remain unchanged and the reversed coefficients in this display decrease. The capped non-one count lies in $[9,31]$. If a non-one count changes, the capped weight is at least 48. If only the ones change, it is at least 42; the only lost midpoint condition is $\bar Y=1^{24}2^9$ at $j=24$. Both compared coefficients then equal $2^9=512$, directly from the append recurrence. In every other case the capped midpoint condition survives.

The finite domain consists of six multiplicities in $[0,24]$, with non-one sum between 9 and 31. All 9,334,875 profiles, and all 158,635,368 applicable comparisons including the exceptional comparison, satisfy~\eqref{eq:mirrorcertificate}; see Section~\ref{sec:certificates}. Finally, a height layer $[l,r]$ above $p_7$ with $l+r<b+5$ would have $2l<b+5$ and $p_l>p_7$. Reflection would put $b+5-l>r$ in the same layer, a contradiction.
\end{proof}

\begin{theorem}\label{thm:central}
Let $4\le M\le7$. Suppose both strict arm weights are at least 19, their difference in absolute value is at most $M-4$, and both arms have at least two non-one parts. If one arm has at least nine non-one parts, then
\[
 \Del H(\min(a,b)+M)>0,\qquad \Del H(\min(a,b)+M+1)>0.
\]
\end{theorem}
\begin{proof}
Orient the long arm as $Y$, put $e=a-b$, and write $D=D_A$, $\delta=D\rev$. At $n=\min(a,b)+M+h$, $h=0,1$, the pure-arm contribution $F_A$ has vanished. A layer $[l,r]$ of $F_Y$ contributes exactly
\[
 d_{n-l}-d_{n-r-1}
\]
to the difference of $F_YD$. If $l+r\ge b+5$, the two indices are ordered and their sum is at most
\[
 2n-b-6\le a+2M+2h-6\le a+M+3,
\]
the kernel reflection sum. Every such contribution is nonnegative.

Let $W$ be the suffix of the last nine non-one parts of $Y$, with forward and reversed coefficients $p_j,q_k$. A word of weight at most nine uses at most nine host positions, so the forward coefficients of $Y$ agree with those of $W$ through rank nine. Lemma~\ref{lem:prefixreverse} gives reverse domination. All 715 possible suffixes over $\{2,\ldots,6\}$ satisfy
\[
 1=q_0\le q_1\le q_2\le q_3\le q_4,
 \qquad q_3\ge p_8,\qquad q_4\ge p_9.
\]
These small finite assertions are included in the suffix check. In particular they hold over each smaller alphabet. Put
\begin{equation}\label{eq:bins}
 m_{jk}=\max\{0,\min(p_j,q_k)-\max(p_{j-1},q_{k-1})\},
 \qquad 2\le j\le9,\quad1\le k\le4.
\end{equation}
The $m_{jk}$ heights start at $j$ and reach a reverse suffix coefficient for the first time at deficit $k$. Their actual right endpoints in $Y$ are at least $b-k$. These bins partition heights $2,\ldots,p_9$.

For $0\le e\le M-4$, define $x_+=\max(0,x)$ and
\begin{align}\label{eq:budgetcentral}
 C_{e,h}={}&\delta_{e-h}-d_{M-1+h}
 +\sum_{j=2}^{7}\sum_{k=1}^{4}m_{jk}(\delta_{j+e-h}-d_{M-1+h+k})\notag\\
 &+\sum_{j=8}^{9}\sum_{k=1}^{4}m_{jk}
                   (\delta_{j+e-h}-d_{M-1+h+k})_+.
\end{align}
Negative reversed indices mean zero. Low kernel coefficients increase through rank eleven, so the right-endpoint estimate gives the displayed coefficient bounds. Layers starting at $j=8$ or 9 are above height $p_7$; their contributions are independently known to be nonnegative. This justifies the positive part. All layers above $p_9$ contribute nonnegatively. Therefore $\Del H(b+M+h)\ge C_{e,h}$. For $e<0$, the reversed indices are those for $C_{0,h}$ and every low index is reduced by $|e|$, so low-rank monotonicity gives the same lower bound.

Suppose first that $A$ has between two and thirteen non-one parts. Cap its multiplicities at thirteen. This preserves $d_0,\ldots,d_{11}$ and decreases $\delta_0,\ldots,\delta_{12}$; every occurrence of a reversed coefficient in~\eqref{eq:budgetcentral} is monotone. The enlarged capped domain has non-one sum between 2 and 13 and either weight at least 19 or a multiplicity equal to 13. Exact exhaustive checks give
\begin{center}
\begin{tabular}{ccrr}\toprule
$M$ & $e$ & $\min C_{e,0}$ & $\min C_{e,1}$\\\midrule
4&0&297&146\\
5&0&291&133\\5&1&298&143\\
6&0&278&112\\6&1&290&130\\6&2&298&142\\
7&0&257&75\\7&1&277&109\\7&2&290&129\\7&3&298&142\\\bottomrule
\end{tabular}
\end{center}

If $A$ has at least fourteen non-one parts, the reversed kernel recurrence gives $\delta_k\ge\binom{14}{k}$. Prepending a one or a two to a word of weight at most seven leaves enough positions in $W$, giving disjoint families. Hence $p_j\ge p_{j-1}+p_{j-2}$ for $2\le j\le8$. With $f_0=f_1=1$, $f_i=f_{i-1}+f_{i-2}$, it follows that $p_j-p_{j-1}\ge f_{j-2}$. Also $p_8\le125$. For $h=0$ all low kernel costs for layers through $p_8$ are at most 253, and for $h=1$ at most 507: use
\[
 d_9\le\sum_{j=2}^{8}2^{j-1}-1=253,\qquad
 d_{10}\le\sum_{j=3}^{9}2^{j-1}-1=507
\]
and low-rank monotonicity. Retaining layers only through height $p_8$ gives the lower bound
\[
 \mathbf1_{e\ge h}\binom{14}{e-h}
 +\sum_{j=2}^{8}f_{j-2}\binom{14}{j+e-h}
 -125\begin{cases}253,&h=0,\\507,&h=1.\end{cases}
\]
For $(e,h)=(0,0),(0,1),\ldots,(3,1)$ these are respectively
\[
 58349,\ 19111,\ 49977,\ 26599,\ 31815,\ 18227,\ 12367,\ 65.
\]
All are positive. This proves the theorem for every non-one count.
\end{proof}

\section{A finite complement from long leading runs}\label{sec:completion}
To cover the complement of Theorem~\ref{thm:central}, a bound on the non-one parts must be accompanied by a valid bound on the leading ones. We give that argument explicitly; no general high-rank single-crossing assertion is assumed.

\begin{lemma}[Two long leading runs]\label{lem:longones}
Let $U,V$ be increasing words with parts in $\{2,\ldots,M-1\}$, allowing empty words. Put $u=\wt U$, $v=\wt V$, $s=\len(U)+\len(V)$ and
\[
 B=M+u+v-s.
\]
If $\alpha,\beta\ge B$, then $F_{1^\alpha U,M,1^\beta V}$ is unimodal.
\end{lemma}
\begin{proof}
Use the polynomials $\RR_U,K_U$ from Lemma~\ref{lem:prefixreverse}, and define
\[
 J_U=(1-z)+(S_M-1)K_U.
\]
Both $K_U$ and $J_U$ are nonnegative: $J_U=K_{U,M}-K_U$ counts the new words beginning with a letter at least two on appending $M$. Their values at one are $K_U(1)=\RR_U(1)$ and $J_U(1)=(M-1)\RR_U(1)$.

The exact padding identity is $(1-z)F_{1^\alpha U}=K_U-z^{\alpha+1}\RR_U$. Substitute it into~\eqref{eq:core}. With $D=M+u+v$, this gives
\begin{equation}\label{eq:padding}
 (1-z)^2H=N_0-z^{\alpha+1}E_1-z^{\beta+1}E_2
                 +z^{\alpha+\beta+2}W,
\end{equation}
where
\begin{align*}
 N_0&=(1-z)(K_U+K_V)+(S_M-1)K_UK_V,\\
 E_1&=\RR_UJ_V,\qquad E_2=\RR_VJ_U,\qquad
 W=(S_M-1)\RR_U\RR_V.
\end{align*}
All four degrees are at most $D$, all four values at one equal $\kappa=(M-1)\RR_U(1)\RR_V(1)$, and $E_1,E_2\ge0$. Since $W$ has no terms below degree $s$,
\[
 D<n<\alpha+\beta+2+s\quad\Longrightarrow\quad\Del^2H(n)\le0.
\]
Although $W$ need not be nonnegative, all its right-tail sums are: for $R=\RR_U\RR_V$ and any integer $k$,
\[
 \sum_{j>k}W_j
 =\sum_{i=1}^M\sum_{h>k-i}R_h-\sum_{h>k}R_h
 \ge(M-1)\sum_{h>k}R_h\ge0.
\]
For $n\ge n_*:=D+\max(\alpha,\beta)+1$, the first three prefix sums in~\eqref{eq:padding} have all finished. Hence $\Del H(n)=-\sum_{j>n-\alpha-\beta-2}W_j\le0$.

The hypotheses imply $\alpha+\beta\ge D$, because $u+v\ge2s$. Lemma~\ref{lem:half} therefore gives nonnegative differences through rank $D$. They also imply $n_*<\alpha+\beta+2+s$. Between $D$ and $n_*$ the differences are nonincreasing; afterward they are nonpositive. There can be no negative difference followed by a positive one.
\end{proof}



\section{A cubic finite-obstruction theorem}\label{sec:cubic}
The quadratic one-arm bound can be combined with a second, independent reduction. A sufficiently large disparity between the non-one skeleton weights forces a single crossing of the central differences. This gives an explicit polynomial bound on the weight of an entire minimal obstruction.

\begin{lemma}[Initial growth of the reversed kernel]\label{lem:reversekernelgrowth}
Let $Y$ be an increasing strict arm and put $e_Y=\wt Y-\len Y$. The reversed kernel $\delta_Y=D_Y\rev$ is nondecreasing through rank
\[
 \left\lfloor\frac{e_Y+M}{2}\right\rfloor.
\]
\end{lemma}
\begin{proof}
We express $\delta_Y$ as a sum of products of intervals, each with center at least $(e_Y+M-1)/2$. The empty case is $[M]$. If $Y=Zt$, with $c=\wt Z$ and $l=\len Z$, use
\[
 \delta_{Zt}=[t]\delta_Z+z^{c+t}[M-t].
\]
Multiplication by $[t]$ raises every center by $(t-1)/2$, exactly matching the increase in $e_Y$. The correction interval has endpoint sum $2c+t+M-1$, whereas the required sum is $c-l+t+M-2$; their difference is $c+l+1>0$. Every interval product is symmetric and unimodal, and hence nondecreasing through the claimed common rank. Sum the inequalities.
\end{proof}

\begin{theorem}[An unbalanced non-one skeleton]\label{thm:skeletonimbalance}
Let $M\ge4$, $X=1^\alpha U$ and $Y=1^\beta V$ be increasing strict arms, with $U,V$ consisting of non-one parts. Put $u=\wt U$, $v=\wt V$, and assume
\[
 u\ge2,\qquad |\wt X-\wt Y|\le M-4,\qquad v\ge4u+3M-12.
\]
Then $Q_{X,Y}$ and $H$ are unimodal.
\end{theorem}
\begin{proof}
Put $a=\wt X$, $b=\wt Y$, $\eta=a-b$, and $R=b+M$. Since $b\ge v\ge4u+3M-12\ge2M-1$, and $V$ has at least two parts, the stronger kernel reflection applies. It gives monotonicity of $D_Y$ through rank $2u+2M-4$. The two-step reflection also makes $Q$ nondecreasing through rank $b+3$, since $\eta\ge4-M$.

Use the nonnegative polynomial $K_U=(1-z)F_U+z\RR_U$ from Lemma~\ref{lem:prefixreverse}, of degree at most $u$. The exact padding identity gives
\begin{equation}\label{eq:skeletondifference}
 \Del Q(n)=(K_UD_Y)_n-(\RR_UD_Y)_{n-\alpha-1}.
\end{equation}
We show that these differences are nonincreasing from rank $b+4$ through rank $R+u+1$.

First, $\len V\le v/2$, so $e_Y=v-\len V\ge v/2\ge2u+M-8$. Lemma~\ref{lem:reversekernelgrowth} therefore makes $\delta_Y$ nondecreasing through rank $L=u+M-4$, including zero extension to negative ranks. Write $K_U^*=z^uK_U(1/z)\ge0$. For $b+4\le n\le R+u$, let $r=u+R-n\le L$. Then
\[
 (K_UD_Y)_n-(K_UD_Y)_{n+1}
 =(K_U^*\delta_Y)_r-(K_U^*\delta_Y)_{r-1}\ge0.
\]
Second, the larger index needed for consecutive coefficients of the negative term in~\eqref{eq:skeletondifference} is at most
\[
 (R+u+1)-\alpha-1=M+2u-\eta\le2u+2M-4.
\]
Convolution with the nonnegative $\RR_U$ preserves the known initial monotonicity of $D_Y$. Thus the subtracted term is nondecreasing throughout this interval. The asserted decrease of $\Del Q$ follows.

For every $n\ge R+u+1$, the first term in~\eqref{eq:skeletondifference} is zero, so $\Del Q(n)\le0$. Before rank $b+4$ the differences are nonnegative, and between these two regions they are nonincreasing. Hence no negative difference can be followed by a positive one. The balanced-band equivalence gives unimodality of $H$.
\end{proof}

\begin{theorem}[Polynomial finite verification bound]\label{thm:cubicbound}
Fix an integer $M\ge4$. If there is a strict core whose relative polynomial $Q$ is not unimodal, then there is one whose full core weight is less than $160M^3$.
Consequently, unimodality of all one-descent compositions of descent peak $M$ can be established by a finite exact check of the strict relative cores below this weight bound.
\end{theorem}
\begin{proof}
Choose a counterexample with the fewest total arm parts. Its arms are nonrectangular and its weights differ by at most $M-4$. By Theorem~\ref{thm:uniformcentral}, we may exchange the arms so that the number of non-one parts in $U$ is less than $N_M=\max\{48,8M(M-3)\}$. Thus
\[
 2\le u\le(M-1)(N_M-1)<8M^2(M-1).
\]
The preceding theorem forces $v<4u+3M-12$. Put $D=M+u+v$ and $s=\len U+\len V$. Then
\[
 D<5u+4M-12<40M^2(M-1)+4M-12<40M^3.
\]
Write $X=1^\alpha U$ and $Y=1^\beta V$. In the balanced band $H$ is also nonunimodal. Lemma~\ref{lem:longones} therefore forces $\min(\alpha,\beta)<D-s$. Weight balance gives
\[
 |\alpha-\beta|\le u+v+M-4=D-4.
\]
Therefore
\[
 \wt{X,M,Y}=D+\alpha+\beta
 \le4D-2s-6<4D<160M^3.
\]
Every reduction used here has an analytic proof. For the finite verification principle in Theorem~\ref{thm:finiteprinciple}, suppose all complete strict cores below the bound have unimodal rank polynomials. A minimal relative counterexample would lie in the balanced band, so its complete core would also be a counterexample below the same bound, a contradiction. Thus all strict relative cores are unimodal. The heavier-arm orientation, repeated-maximum recurrence, and restoration of the increasing tail then give the one-descent statement as in the proof of Theorem~\ref{thm:main}.
\end{proof}

\begin{corollary}[Linear bound for a single non-one part]\label{cor:linearshort}
If a minimal counterexample at peak $M\ge4$ has an arm with exactly one non-one part, then its total core weight is less than $36M$.
\end{corollary}
\begin{proof}
Orient that arm as $X$. Then $u\le M-1$, and Theorem~\ref{thm:skeletonimbalance} forces $v<4u+3M-12$. Thus $D=M+u+v<9M-17$. The same long-run and balance calculation gives total weight less than $4D<36M$.
\end{proof}

The finite verification bound is polynomial in the peak; the number of cores below it need not be polynomial. No exhaustive computation of this general domain has been performed here, and no unrestricted unimodality conclusion is inferred. This bound is for relative unimodality; the relative tail-log-concavity bound reported in Chen's abstract concerns a stronger property.

\section{An arm with one non-one part}\label{sec:single}
\begin{lemma}\label{lem:single}
Let $4\le M\le7$, $X=1^\alpha t$ with $2\le t<M$, and $|\wt X-\wt Y|\le M-4$. If $Y$ has at least nine non-one parts, then $Q_{X,Y}$ is unimodal.
\end{lemma}
\begin{proof}
Put $b=\wt Y$, $\eta=\wt X-b$, $R=b+M$, and $q_i=(F_Y)_{b-i}$. The cases $b<35$ are the direct finite bases in Section~\ref{sec:certificates}. Assume $b\ge35$. With $D=D_Y$ and $P=S_t[\alpha+1]D$, we have $Q=D+P$ and $P$ is unimodal. The kernel grows through rank 21, while two-step reflection gives growth of $Q$ through at least $R-4$. The exact window formula is
\begin{equation}\label{eq:singlewindow}
 \Del P(n)=\sum_{j=1}^t(d_{n-j}-d_{n-\alpha-j-1}).
\end{equation}

For $c=0,1,2,3$, telescoping gives
\begin{align}\label{eq:singlegaps}
 \Del Q(R-c)-\Del Q(R-c+1)
 ={}&q_c+\sum_{j=c+2}^{c+t}q_j-T_{c+t-M}\notag\\
 &+d_{M-c+t-\eta-1}-d_{M-c-\eta-1}
   +2\mathbf1_{M=4,\ c=3},
\end{align}
where $T_s=0$ for $s<0$, and $T_s=\sum_{j=0}^{s-1}q_j+2q_s$ for $s\ge0$. To check the endpoint correction, $d_{R-x}=\sum_{j=0}^xq_j$ for $x<M$; when $x=M+s$, the latter sum must be reduced by $T_s$. Only $s=0,1,2$ occurs. The extra two arises from $d_{R-c-1}=d_b$ when $M=4,c=3$. The displayed low kernel difference is nonnegative, since its indices lie between zero and fifteen.

Let $N\ge9$ be the number of non-one parts in $Y$. Its reversed recurrence gives
\[
 q_0=1,\quad q_1\le N+1,\quad q_2\le\binom{N+2}{2},
 \qquad q_j\ge\binom Nj.
\]
For the upper bounds, an initial run of ones supplies at most one extra geometric factor. A correction at degree two can occur only with no initial ones and first part two, still within the displayed bound. Thus $T_s\le(N+2)^2$ for $s\le2$. Both $\binom N4$ and $\binom N5$ are at least $(N+2)^2$ for $N\ge9$: check $N=9$, and compare successive ratios. If a correction occurs and $c=3$, the positive sum contains $q_5$ and pays for it. If $c=2$, it contains $q_4$. If $c=1$, only $T_0=2$ occurs, paid by $q_3$; for $c=0$ no correction occurs. In all cases $q_c>0$ remains. The five differences at $R-3,\ldots,R+1$ therefore strictly decrease.

For $R+1\le n\le R+t$, another telescoping of~\eqref{eq:singlewindow} gives
\[
 \Del P(n+1)-\Del P(n)
 =-d_{n-t}-(d_{n-\alpha-1}-d_{n-\alpha-t-1})\le0,
\]
because $d_n=0$ and the low indices are at most 21. For $n\ge R+t+1$, all high terms in~\eqref{eq:singlewindow} vanish, so $\Del P(n)\le0$. If any of the five central differences of $Q$ is negative, their strict decrease and integrality give $\Del P(R+1)=\Del Q(R+1)+1\le0$, and no later increase occurs. Otherwise the increasing segment joins the unimodal tail $P$. This proves the lemma.
\end{proof}

\section{Completion of the bounded-peak theorem}
\begin{proposition}\label{prop:finitecomplement}
Let $4\le M\le7$. If both strict arms have at most eight non-one parts and their weights differ by at most $M-4$, their complete core polynomial is unimodal.
\end{proposition}
\begin{proof}
Orient the weights with difference $d\in\{0,\ldots,M-4\}$ and write $X=1^\alpha U$, $Y=1^\beta V$, with $u=\wt U$, $v=\wt V$, $s=\len U+\len V$. Then
\[
 \beta=\alpha+u-v-d,\quad \alpha_0=\max(0,v+d-u),\quad B=M+u+v-s.
\]
The minimum of the two one-run lengths is $\alpha-\alpha_0$. Lemma~\ref{lem:longones} covers all $\alpha\ge\alpha_0+B$. The remaining finite domain is exactly
\begin{equation}\label{eq:paddingdomain}
 \begin{gathered}
 U,V\text{ increasing over }\{2,\ldots,M-1\},\quad \len U,\len V\le8,\\
 0\le d\le M-4,\quad \alpha_0\le\alpha<\alpha_0+B,\quad
 \beta=\alpha+u-v-d.
 \end{gathered}
\end{equation}
Complete exact computations establish unimodality of all these polynomials. The counts for $M=4,5,6,7$ are respectively $40,500$, $1,579,050$, $27,932,850$, and $311,397,372$. Two independently organized implementations, described in Section~\ref{sec:certificates}, agree throughout. The long-run lemma covers every omitted run length.
\end{proof}

\begin{proof}[Proof of Theorem~\ref{thm:main}]
First prove unimodality of every strict $Q$ for each $M\le7$. Peaks one and two are immediate, and Proposition~\ref{prop:outer} excludes a minimal counterexample at peak three. For $4\le M\le7$, orient a minimal counterexample as $a=b+d$. Proposition~\ref{prop:outer} gives $0\le d\le M-4$. The complete small bases exclude $b\le18$, so $b\ge19$. Rectangular arms are excluded by Lemma~\ref{lem:balanced}. Proposition~\ref{prop:finitecomplement} handles the case when both non-one counts are at most eight. Otherwise one arm has at least nine non-one parts. Lemma~\ref{lem:single} excludes a single non-one part in the other arm. Theorem~\ref{thm:central} gives positive differences at $R=b+M$ and $R+1$.

If $M\le6$, or $M=7$ and $d\ge1$, the stronger kernel reflection gives growth through
\[
 \left\lfloor\frac{a+b+M+4}{2}\right\rfloor\ge R-1.
\]
If $M=7,d=0$, choose $Y$ to have at least nine non-one parts. Each higher layer of $F_X$ has endpoint sum at least $b+1$; multiplying it by $D_Y$ gives growth through $R-1$. The lowest layer contributes $d_{R-1}-d_5\ge10-8>0$. The first bound uses the nine non-one parts in the reversed coefficient; the second uses $p_0+\cdots+p_4\le16$ and $p_5\ge8$, since every composition of five with parts one or two embeds in $Y$. Thus $Q$ grows through $R-1$ in this case too.

The two strict central differences extend growth through $R+1$. Delete the last part $u$ of the heavier arm. The shorter $Q_{X^-,Y}$ is unimodal by minimality. At every rank at least $R+1$, the recurrence $Q=D_Y+S_uQ_{X^-,Y}$ agrees with the unimodal second summand. The increasing segment therefore joins a unimodal tail, a contradiction. By Lemma~\ref{lem:balanced}, orienting any strict core with the heavier arm first now gives unimodality of $H$.

For a general one-descent host whose descending upper part is $M\le7$, write
\[
 w=(U,M^{r+1},Y,E),
\]
where $U,Y$ are increasing with parts strictly below $M$, and $E$ is increasing with parts at least $M$. Start with the strict core $(U,M,Y)$ and restore each extra central maximum using
\[
 F_{U,M^{r+2},Y}=1+S_MF_{U,M^{r+1},Y}.
\]
This follows directly from the maximum decomposition. Then append each part of $E$, a current global maximum. Every operation is $P\mapsto1+S_tP$ and preserves unimodality for these polynomials with constant term one. Monotone hosts are already covered by Lemma~\ref{lem:arms}.

Finally, reversal of smaller words preserves weight, so $F_Y=F_{Y^{\mathrm{rev}}}$. The maximum decomposition makes $(X,M,Y)$ equiranked with $(X,M,Y^{\mathrm{rev}})$. This covers every bitonic host of maximum at most seven, including repeated maxima.
\end{proof}

\section{Scope and further questions}
The proof covers arbitrary length at peaks through seven. It does not prove the unrestricted composition conjecture or tail log-concavity. At arbitrary peaks, Theorem~\ref{thm:finiteprinciple} reduces the problem to a specified finite domain with a cubic weight cutoff. The number of candidates may still be very large. The profile certificates through peak seven use substantially smaller specialized domains; no complete peak-eight computation is claimed. A failed lower bound need not be a failed rank inequality: the positive-part refinement uses structural information discarded by the unrefined estimate. Further progress could come from sharper finite reductions or uniform positivity within the remaining bounded domains. No conclusion at further peaks is inferred from the present computations.


The effective dilation theorem gives a different finite reduction:
for each fixed shape, only finitely many dilations remain below its
certified threshold. It does not provide a threshold uniform in all
shapes. The lower eventual mode and the mean have the same center
$(tW+l-1)/2$ up to rounding and a vanishing mean error, respectively.
Sharper bounds for the dilation threshold, and extensions of the
cumulative theorem to arbitrary hosts, are natural further questions.

\section{Exact certificates and reproducibility}\label{sec:certificates}
Every finite domain used above is explicit. Each program returns failure on a violated assertion; successful completion is reported only after the full domain finishes. Random diagnostics are separate from the proof inputs.

\subsection{Forward and reversed profiles}
Starting at weight zero with $f_0=q_0=1$ and other entries zero, append $t$ using
\begin{equation}\label{eq:profilealgorithm}
 \begin{split}
 f'_0&=1,\qquad f'_k=\sum_{j=1}^{\min(t,k)}f_{k-j}\quad(k>0),\\
 q'_k&=\sum_{j=0}^{\min(t-1,k)}q_{k-j}+\mathbf1_{k=\mathrm{weight}+t}.
 \end{split}
\end{equation}
The mirror certificate uses forward degrees through 24 and reversed deficits through 19. It enumerates $c_1,\ldots,c_6\in[0,24]$ with $9\le c_2+\cdots+c_6\le31$, checks~\eqref{eq:mirrorcertificate}, and includes the exceptional capped profile specified there. The profile count is
\[
 25\left(\binom{36}{5}-5\binom{11}{5}-\binom{13}{5}\right)=9,334,875.
\]
There are exactly 158,635,368 applicable coefficient comparisons. Both implementations pass all of them. The C program \texttt{mirror.c} saturates reversed entries at $2^{23}$. This is exact for these comparisons: every tested forward coefficient is at most $2^{23}$, and saturation of a nonnegative sum commutes with saturation of its summands. The independent Python mirror checker uses arbitrary-precision integers without saturation.

For the central certificate, compute
\[
 d_k=\mathbf1_{k=0}-f_k+\sum_{j=1}^Mf_{k-j},\qquad
 \delta_k=\sum_{j=0}^{M-1}q_{k-j}-q_{k-M}.
\]
The omitted reversal term at degree $\wt A+M$ is beyond every deficit used. The short profile has multiplicities in $[0,13]$, non-one sum in $[2,13]$, and either weight at least 19 or some multiplicity equal to 13. The suffix has nine parts in $\{2,\ldots,M-1\}$. Each pair is tested for $0\le e\le M-4$ and $h=0,1$.
\begin{center}
\begin{tabular}{crrrr}\toprule
$M$ & Short profiles & Suffixes & Pairs & Inequalities per pair\\\midrule
4&1,209&10&12,090&2\\
5&7,339&55&403,645&4\\
6&32,565&220&7,164,300&6\\
7&118,969&715&85,062,835&8\\\bottomrule
\end{tabular}
\end{center}
\texttt{central\_all.c} evaluates~\eqref{eq:budgetcentral} directly. The independent \texttt{independent\_profiles.py} constructs matrices for the low-layer coefficients and separate positive-part gains for the higher layers. It also checks the suffix endpoint assertions. Both reproduce the minima in Theorem~\ref{thm:central}.

All C and NumPy central arithmetic is safe in signed 64-bit integers. A short profile has length at most 26; each reversed coefficient through deficit twelve is at most $27\binom{37}{12}$. There are at most 248 relevant nonlowest heights. An unconditional bound for every intermediate absolute sum is
\[
 249\cdot8\cdot27\binom{37}{12}+249\cdot16384
 =99,633,949,536,480<2^{63}.
\]
The analytic bounds for fourteen or more non-one parts are also evaluated separately with Python integers.

\subsection{Small bases and the bounded complement}
\texttt{small\_bases.py} constructs rank polynomials directly from the append recurrence and the maximum decomposition. The balanced domain is $1\le b\le18$, $0\le d\le M-4$, with all increasing arms of weights $b+d,b$. The single-arm domain is $1\le b<35$, all increasing $Y$ of weight $b$ with at least nine non-one parts, $|\eta|\le M-4$, and $X=1^{b+\eta-t}t$ for $2\le t<M$ whenever the exponent is nonnegative.
\begin{center}
\begin{tabular}{crrrr}\toprule
$M$ & Balanced bases & Single-arm bases & Non-one tails & Padded cores\\\midrule
4&6,329&972&45&40,500\\
5&56,521&12,420&165&1,579,050\\
6&234,073&51,500&495&27,932,850\\
7&639,460&131,110&1,287&311,397,372\\\bottomrule
\end{tabular}
\end{center}
The direct Python bases use base-$2^{64}$ encoding for exact convolution. No slot can carry: balanced hosts have total weight at most 46, while in the single-arm bases $F_X(1)\le229$, $F_Y(1)\le2^{34}$, and $Q(1)\le229(1+6\cdot2^{34})<2^{64}$.

The final column is the complete domain~\eqref{eq:paddingdomain}. If $T=\binom{M+6}{M-2}$ is the number of non-one tails of length at most eight, the number of padded instances is
\[
 (M-3)T\left(MT+2\sum_U(\wt U-\len U)\right)
 =(M-3)(9M-16)T^2.
\]
The second equality follows by viewing each tail multiplicity vector, together with its unused slots, as a weak composition of eight into $M-1$ coordinates. Symmetry gives $\sum_U(\wt U-\len U)=4(M-2)T$.
Every host in the domain has weight at most 278: with $r,t\le8$, $u\le6r$, $v\le6t$, and $0\le d\le3$,
\[
 \wt{X,M,Y}\le3(u+v)-2(r+t)+19+|u-v-d|\le278.
\]
The one-run lengths are less than 145, and full rank-polynomial masses are below $2(145\cdot6^8)+7(145\cdot6^8)^2<2^{59}$.

\texttt{bounded\_tails.c} evaluates the cumulative second-difference formula~\eqref{eq:padding}. Independently, \texttt{direct\_padding.c} initializes a full rank polynomial by ordinary convolution, then maintains its coefficients. When the two one-runs increase by one, put $x=z^{\alpha+1}\RR_U$, $y=z^{\beta+1}\RR_V$. The update is
\[
 H(A+x,B+y)-H(A,B)=xD_B+yD_A+(S_M-1)xy.
\]
The direct implementation checks coefficient nonnegativity and unimodality after every update. Both programs complete all four domains with zero failures. Their intermediate absolute values are bounded by $2^{47}$ for cumulative second differences and $2^{60}$ for the direct update. The largest actual full coefficient, at peak seven, is $1,403,945,841,784,095$.

\subsection{Lean scope and trust boundary}
The eleven supplied Lean files were checked using Lean 4.34.0 with the installed mathlib environment. None contains a \texttt{sorry} or \texttt{admit} declaration. The printed axiom dependencies are included with the results.

\texttt{SevenMirror.lean} defines the complete reduced mirror enumeration, including the exceptional profile, and proves its Boolean truth by \texttt{native\_decide}. It also checks the nineteen binomial inequalities at $N=32$. These theorems explicitly depend on native-evaluation axioms; they are not represented as pure kernel reduction proofs of the large computation.

\texttt{LayerBridge.lean} contains ordinary kernel proofs of the passage from reflection to layer centers, positive layer budgets, and the algebraic identities for prefix reversal and repeated maxima. \texttt{SevenBridge.lean} checks the high-layer index inequality for $M\le7$, clipping and its monotonicity under capping, the weight-278 bound, the direct-padding update, single-arm payment inequalities, and auxiliary uniform index and budget steps. \texttt{QuadraticBridge.lean} proves the unbounded block-binomial inequality $t^k\le\binom{tk}{k}$, the central deficit range, the closed quadratic cutoff formula, the strict power comparison, and the finite-sum statement that one positive contribution pays all low-layer costs in Theorem~\ref{thm:uniformcentral}. \texttt{CubicBridge.lean} checks the skeleton-disparity hypotheses, both central index ranges, the single-crossing implication, the cubic weight bound, and the linear bound for a single non-one part. \texttt{DeficitGrowth.lean} proves coefficientwise differential inequalities for the actual unbounded reversed append recurrences. Its input list stores host parts from right to left; the identification of these recurrences with distinct-word rank polynomials remains a written combinatorial proof. These use only ordinary logical axioms recorded in the logs. The combinatorial interpretation, kernel shape lemmas, central matrix certificate, padding enumerations, and complete main induction have not been fully formalized in Lean.

Five further files concern dilation, cumulative regularity, sharp
growth, and a dominant maximum.

\texttt{CumulativeReverse.lean} proves log-concavity for the actual
window-plus-terminal-step recurrence on arbitrary positive lists; this
includes the window-convolution proof and both boundary comparisons.
\texttt{BinarySaturation.lean} proves that a nonempty fixed-length binary
language closed under both shifts is universal.
\texttt{EscapeDeficit.lean} proves the bound
$|S|\le(2|S|-|T|)(2^h-1)$ for an arbitrary finite set of binary states
$S$, a set of feasible directed edges $T$, and a prescribed block of
$h$ shifts forcing every state to leave $S$. It verifies the two-preimage
bound and the comparison between outgoing edges and the rank deficit.
Its final theorem derives the forced escape condition from closure
under contiguous factors and a missing binary word, for arbitrary
lengths; the condition is not left as an additional hypothesis there.
These are ordinary kernel proofs with only the standard logical axioms.
The translation from compositions to cut strings and the cited
cumulative-convolution closure theorem remain written proof steps.
\texttt{DominantMaximum.lean} proves the central product inequality
from its endpoint budgets and the unbounded sign argument excluding a
decrease followed by an increase when $N\le2M+5$. The endpoint counting
and Boolean-lattice inputs are proved in the text, not formalized by
that file.
\texttt{DilationBridge.lean} proves the unbounded coefficient-dominance
bound for real polynomials, its ratio-bound corollary used in $b(P)$,
and the passage from the three rank regions to a single crossing.
The multilinear composition expansion, central curvature, and
identification of the exact central polynomials remain written
mathematical proofs.

\subsection{Reproduction and provenance}
Direct enumeration checks the growth bounds on 1,112 downsets, using
9,973 rank comparisons. The shift-graph estimate is checked on 65,808
subsets. Separate cumulative diagnostics check 6,434 exhaustive and
4,000 random monotone arms, comprising
3,640,541 coefficient positions. A further 16,000-core diagnostic checks
20,438,112 cumulative comparisons, including restored terminal maxima.
A separate rightmost-last-letter oracle is checked against direct
distinct-word enumeration on all 512 hosts of weight at most nine.
It then checks the valley rearrangement, common-maximum factorization,
and 5,631,724 cumulative comparisons on 6,000 hosts, with up to 36 observed descents,
including subsequent larger terminal maxima.
A separate dominant-maximum diagnostic checks all 8,191 nonempty
hosts of weight at most thirteen for the endpoint inequalities, and
6,256 arm pairs producing 31,280 dominant-maximum hosts.
The dilation diagnostic checks 70,113 exact coefficient differences,
9,560 central-polynomial evaluations, 4,780 signs at million-scale
dilations, and 2,460 evaluations of the exact mode-choice polynomial.
Independently, 1,213 distinct balanced shapes are replayed at their
analytically certified thresholds: 2,426 full rank sequences and
8,678 rational-polynomial evaluations pass. The largest sufficient
threshold in this diagnostic is 672; it is not a universal bound.
The right-tail bound is also checked on 26,383 hosts, using 658,801
coefficient comparisons.
These finite diagnostics test formulas; the unbounded statements follow
from the analytic proofs.

With Python 3, NumPy, and a C compiler, run the supplied \texttt{Run-Certificates.ps1}, or run the individual programs. The mirror program takes arguments \texttt{24 9 7}; the central C and Python programs and the direct-padding program accept a peak from four through seven. The cumulative-padding program checks all four peaks. Python profiles and small bases are independent of the C path.

Diagnostics in \texttt{audit\_unified.py} and \texttt{audit\_quadratic.py} compare the formulas with actual polynomials. The latter checks 27 arm pairs at peaks four through twelve and all 97 central positions in the theorem's range. The file \texttt{audit\_uniform.py} records an earlier diagnostic of the larger-count argument. The cubic reduction has an additional 375-core exact diagnostic, including 10,423 reversed-kernel growth comparisons and 18,317 central monotonicity positions. The reversed append estimates are also checked on 19,448 arms. Diagnostics do not replace analytic proofs or exhaustive certificates.

The Lean runner accepts the local Lean executable and a mathlib package directory. The manuscript compiles from \texttt{paper.tex} using a standard LaTeX distribution or Tectonic. No network is required by any arithmetic checker. The manifest records the exact source and result files distributed with the paper.

The maximum decomposition, elementary arm facts, and interval methods
are background tools. The required versions, kernel estimates, peeling
steps, and long-run lemma are proved in the manuscript. In particular,
the padding domain follows from the proved long-run lemma and does not
assume a general likelihood-ratio comparison for uncumulated kernels.

The finite-profile argument closes peaks six and seven. Separately,
Theorem~\ref{thm:uniformcentral} gives a quadratic one-arm bound and
Theorem~\ref{thm:cubicbound} gives a cubic total-weight bound at every
peak. The other main results establish cumulative regularity without a
peak bound and sharp growth deficits for arbitrary composition
downsets. Exact arithmetic and the stated Lean checks establish the
supplied certificates. No claim about an inaccessible paper's full
contents, or about the complete formalization of the manuscript, is
made.

The dilation certificate program accepts a list of positive parts.
It constructs exact central polynomials and their coefficient bounds,
returning a threshold that applies to every subsequent integer
dilation. The supplied example for $(1,4,2,5,3,2,4)$ has threshold 32.
Full-profile diagnostic replays at that threshold are separate from
the analytic proof that the certificate covers all larger dilations.


\clearpage
\begin{thebibliography}{9}
\raggedright
\bibitem{AlbertEtAl} M. Albert, M. Engen, J. Pantone and V. Vatter, \emph{Universal layered permutations}, Electronic Journal of Combinatorics \textbf{25} (2018), P3.23. \url{https://arxiv.org/abs/1710.04240}.
\bibitem{NandaSengupta} A. K. Nanda and D. Sengupta, \emph{Discrete Life Distributions with Decreasing Reversed Hazard}, Sankhy\={a}: The Indian Journal of Statistics \textbf{67} (2005), 106--125. \url{https://www.jstor.org/stable/25053410}.
\bibitem{Sagan} B. E. Sagan, \emph{Compositions inside a rectangle and unimodality}, Journal of Algebraic Combinatorics \textbf{29} (2009), 405--411. \url{https://arxiv.org/abs/0707.1052}.
\bibitem{Vatter} V. Vatter, \emph{An assortment of problems in permutation patterns: unimodality, equivalence, derangements, and sorting}, arXiv:2602.16355v2. \url{https://arxiv.org/abs/2602.16355}.
\bibitem{VatterWords} V. Vatter, \emph{Log-concavity of subsequence counts of words}, arXiv:2608.22147. \url{https://arxiv.org/abs/2608.22147}.
\bibitem{MenonSingh} K. Menon and A. Singh, \emph{Pattern avoidance and dominating compositions}, arXiv:2104.07274. \url{https://arxiv.org/abs/2104.07274}.
\bibitem{Zhang} Z. Zhang, \emph{Rank-Unimodality for One-Descent Families in Generalized Subword Order}, SSRN 7397439, posted 5 September 2026. \url{https://ssrn.com/abstract=7397439}. Public abstract and the matching anonymized manuscript in the supplied archive consulted; identity of the complete public version not independently confirmed.
\bibitem{Chen} L. Chen, \emph{Rank-unimodality and tail log-concavity of composition ideals}, SSRN 7385218, revised 6 October 2026. \url{https://ssrn.com/abstract=7385218}. Abstract consulted; full text not independently retrieved in this work.
\end{thebibliography}
\end{document}








